% Experiments-section fragment for the SCJEPA manuscript.
% Intended to be included from the main paper with % Experiments-section fragment for the SCJEPA manuscript.
% Intended to be included from the main paper with % Experiments-section fragment for the SCJEPA manuscript.
% Intended to be included from the main paper with % Experiments-section fragment for the SCJEPA manuscript.
% Intended to be included from the main paper with \input{docs/experiments}.
% Required packages: amsmath, amssymb, bm, booktabs.
% Status: Experiment 1 is implemented and running; Experiments 2 and 3 are
% prospective protocols. Replace the explicit status statements with measured
% results once each continuation gate has been passed.

\section{Experiments}
\label{sec:experiments}

We evaluate whether sparse predictive structure can identify time-invariant physical
parameters while progressively removing access to the ground-truth state.  We use a
three-experiment ladder.  Experiment~1 gives both branches the true physical state;
Experiment~2 replaces the predictor-side state by a causal visual encoding while retaining a
true-state target; and Experiment~3 jointly learns the context and single-frame target
representations from pixels.  The physical system, parameter distribution, SPARTAN predictor,
and evaluation logic are held fixed as far as the change in observation contract permits.
Each experiment is run only after its predecessor passes a predeclared continuation gate.  This
staging separates failures of causal parameter identification from failures of perception,
tracking, or jointly learned target geometry.

At the time of writing, Experiment~1 is implemented and its attention-logit sweep is running.
The later experiments below specify the intended confirmatory protocols and must not be read as
completed empirical results.

\subsection{Bounce}
\label{sec:experiments:bounce}

\subsubsection{Data}
\label{sec:experiments:bounce:data}

An episode contains $N=5$ elastic balls in the unit square.  We use $X_t$ for a rendered
frame, $Z_t=(z_t^1,\ldots,z_t^N)$ for the true kinematic state, and
$\boldsymbol{m}=(m_1,\ldots,m_N)$ for the time-invariant physical parameters, where
\begin{equation}
    X_t\in[0,1]^{3\times64\times64},
    \qquad
    z_t^i=(p_{x,t}^i,p_{y,t}^i,v_{x,t}^i,v_{y,t}^i)\in\mathbb{R}^{4},
    \qquad
    m_i>0.
    \label{eq:bounce-observations}
\end{equation}
For every episode and object, the implemented Baumgartner-aligned configuration samples
\begin{equation}
    \epsilon_i\sim\mathcal{N}(0,1),
    \qquad
    m_i=\operatorname{clip}_{[0.5,3.0]}(1.5+0.5\epsilon_i),
    \qquad
    r_i=r_0\frac{m_i}{m_{\mathrm{ref}}},
    \label{eq:bounce-mass-radius}
\end{equation}
with $r_0=0.08$ and the episode-independent reference $m_{\mathrm{ref}}=1.5$.
Equation~\eqref{eq:bounce-mass-radius} is a clipped Gaussian, rather than a truncated
Gaussian, and therefore places atoms at the clipping boundaries.  The numerical mass
distribution is our declared choice because its exact parameters are not specified by the
source Bounce experiment.

The initial velocity has fixed norm $0.7$ and uniformly random direction,
\begin{equation}
    \varphi_i\sim\mathcal{U}[0,2\pi),
    \qquad
    v_0^i=0.7(\cos\varphi_i,\sin\varphi_i).
\end{equation}
The current generator sequentially samples an initial centre uniformly from
$[1.05r_i,1-1.05r_i]^2$ and accepts it only if
$\lVert p_0^i-p_0^j\rVert_2>1.1(r_i+r_j)$ for all previously placed objects $j$.
Consequently, the support of the initial positions is itself mass-dependent.  We report this
fact rather than describing the initial state as independent of mass.  To isolate evidence from
the subsequent dynamics, the final study also includes a mass-independent-initialization
control: all wall margins and pair-separation constraints are computed using the common
worst-case radius $r_{\max}=r_0(3.0/1.5)$, irrespective of the sampled masses.  This control is
required before attributing recovery specifically to transition dynamics rather than to the
initial-state support.

The simulator records $T=60$ states at intervals $\Delta t=0.1$, with eight physics
substeps per recorded transition.  Free motion is followed by exact reflection of wall
overshoot about $r_i$ or $1-r_i$.  For a colliding pair, let
\begin{equation}
    n_{ij}=\frac{p_i-p_j}{\lVert p_i-p_j\rVert_2},
    \qquad
    u_{ij}=(v_i-v_j)^\top n_{ij}<0.
\end{equation}
The restitution-one velocity update is
\begin{align}
    v_i' &= v_i-\frac{2m_j}{m_i+m_j}u_{ij}n_{ij},
    &
    v_j' &= v_j+\frac{2m_i}{m_i+m_j}u_{ij}n_{ij}.
    \label{eq:elastic-collision}
\end{align}
Before applying this impulse, any penetration is projected back to the contact surface
$\lVert p_i-p_j\rVert_2=r_i+r_j$ using inverse-mass weights.  The corrected simulator is
versioned so that pre-generated trajectories from older collision rules cannot be mixed with
the current data.

Each episode supplies
\begin{equation}
    (X_{0:T-1},Z_{0:T-1},\boldsymbol{m},C_{0:T-2}),
    \qquad
    C_t\in\{0,1\}^{N\times N},
\end{equation}
where $C_{t,ij}=1$ records a pair contact during transition $t\rightarrow t+1$.
Off-diagonal entries are symmetric.  In the mass-dependent-radius environment,
$C_{t,ii}=1$ records a wall contact because its location depends on $m_i$ through $r_i$.
We use the first $T_{\mathrm{par}}=30$ observations to infer one parameter representation and
evaluate $K=T-T_{\mathrm{par}}=30$ transitions,
\begin{equation}
    \mathcal{I}=\{T_{\mathrm{par}}-1,\ldots,T-2\}=\{29,\ldots,58\}.
    \label{eq:transition-index-set}
\end{equation}
Thus, the parameter encoder sees observations $0,\ldots,29$, while the supervised pairs in
Experiment~1 are $(Z_{29},Z_{30}),\ldots,(Z_{58},Z_{59})$.  These are 30
teacher-forced one-step predictions, not one 30-step open-loop rollout.

\paragraph{Why absolute masses are observable.}
If all physical radii were equal, multiplying every mass by a common positive constant would
leave Equation~\eqref{eq:elastic-collision} unchanged; elastic impulses would then identify
mass ratios but not absolute scale.  The implemented relation
$r_i=r_0m_i/m_{\mathrm{ref}}$ removes that symmetry because a common mass rescaling also
changes pair-contact distances, wall-contact locations, penetration correction, and the support
of initial positions.  The current state experiment sets rendering off and never supplies a
literal radius coordinate, but its observed positions and velocities are nevertheless generated
by this mass-dependent geometry.  We therefore claim identification in the
mass-proportional-physical-radius environment, not in the harder equal-radius environment.

For Experiments~2 and~3, physical and rendered radii must be separated.  The simulator retains
Equation~\eqref{eq:bounce-mass-radius}, while every object is drawn with the same
mass-independent glyph and rendered radius.  The primary visual data must also contain no
dataset-global visual signature tied to simulator row index.  Hence a single frame cannot reveal
mass from the object's glyph, although a sequence can reveal mass through centre motion and
hidden contact geometry.  A visible-radius version and an equal-physical-radius version are
reported only as labelled controls.  The current data implementation uses a single radius switch
for physics and rendering; splitting those two settings and versioning the resulting data are
prerequisites for the visual experiments.

\subsubsection{Ground-truth local graphs}
\label{sec:experiments:bounce:graphs}

The destination row is indexed first in every graph.  From $C_t$, we define the true
state-to-state and mass-to-state local graphs as
\begin{align}
    G^{Z\rightarrow Z}_{t,ij}
        &= \mathbf{1}\{i=j\ \lor\ C_{t,ij}=1\},
        \\
    G^{m\rightarrow Z}_{t,ij}
        &= \mathbf{1}\!\left\{C_{t,ij}=1\ \lor\
             \left(i=j\ \land\ \bigvee_k C_{t,ik}=1\right)\right\}.
    \label{eq:bounce-ground-truth-graphs}
\end{align}
A pair collision therefore gives each affected state an own-mass and partner-mass parent; a
mass-dependent wall collision gives the affected state an own-mass parent.  Self-state edges are
always present because free motion depends on the object's previous state.

\subsubsection{Shared model comparisons and metrics}
\label{sec:experiments:shared-protocol}

Every scientific comparison uses the same three predictor modes under matched initialization
rules, training budgets, and data splits:
\begin{enumerate}
    \item \emph{Dense}: $A^{(\ell)}\equiv\boldsymbol{1}$ in every SPARTAN layer and no
          path penalty.
    \item \emph{Token-local}: $A^{(\ell)}\equiv\boldsymbol{0}$, so only each token's
          residual--MLP path remains.  This reference has no access to other state or parameter
          tokens, but its own state may contain statistical traces of earlier mass-dependent
          geometry; it is therefore more precise than calling the reference fully mass-blind.
    \item \emph{Sparse}: learned hard Bernoulli gates with the SPARTAN path objective and a
          GECO-style constraint controller.
\end{enumerate}

There are $N$ state tokens and $N$ parameter tokens.  For $L_{\mathrm{sp}}=3$ SPARTAN
layers, the multilayer reachability matrix is
\begin{equation}
    \overline A
    =\bigl(A^{(L_{\mathrm{sp}})}+I\bigr)\cdots
     \bigl(A^{(1)}+I\bigr),
    \label{eq:spartan-path-matrix}
\end{equation}
where $\overline A_{ij}>0$ means that at least one permitted computational path connects source
token $j$ to decoded output $i$.  The optimized path count includes only the $N$ decoded state
rows,
\begin{equation}
    \mathcal{L}_{\mathrm{path}}
    =\mathbb{E}\!\left[
       \sum_{i=1}^{N}\sum_{j=1}^{2N}\overline A_{ij}
     \right],
    \qquad
    \rho_{\mathrm{path}}
    =\frac{1}{2N^2}\sum_{i=1}^{N}\sum_{j=1}^{2N}
       \mathbf{1}\{\overline A_{ij}>0\}.
    \label{eq:path-objective-density}
\end{equation}
The second quantity discards path multiplicity and is the primary pruning metric.  With ten
tokens, the token-local reference has $\overline A=I_{10}$ and
$\rho_{\mathrm{path}}=5/50=0.1$, while the dense model has
$\rho_{\mathrm{path}}=1$.  The corresponding unnormalised path objectives are 5 and 6655:
if $J$ is the $10\times10$ all-ones matrix, then
$(J+I)^3=I+133J$, each decoded row sums to 1331, and five decoded rows sum to 6655.
Reachability over all output rows is retained as a diagnostic only, since parameter-token outputs
are not decoded.

For a fixed target space, training minimizes
\begin{equation}
    \mathcal{L}
    =\mathcal{L}_{\mathrm{pred}}
     +\lambda_{\mathrm{logit}}\mathcal{L}_{\mathrm{logit}}
     +\lambda_{\mathrm{reg}}\mathcal{R}
     +\Lambda^{-1}\mathcal{L}_{\mathrm{path}},
    \label{eq:shared-training-objective}
\end{equation}
when sparsity is enabled.  Here $\mathcal{R}$ is VISReg on learned
representations.  Fixed raw-state targets are not regularized: applying VISReg to them produces
no model gradient and only adds a stochastic constant.  The attention-logit penalty averages
$\exp(a)+\exp(-a)$ over layer logits $a$ and has theoretical minimum 2; the implementation uses
a linear continuation beyond $|a|=30$ for numerical stability.

The dual constraint excludes both VISReg and the path penalty,
\begin{equation}
    c=\mathcal{L}_{\mathrm{pred}}
      +\lambda_{\mathrm{logit}}\mathcal{L}_{\mathrm{logit}}
      \leq \tau.
    \label{eq:raw-geco-constraint}
\end{equation}
For every architecture and seed, $\tau$ is recalibrated as the held-out constraint of a
converged dense reference, and the sparse run is launched only if the dense constraint is below
the token-local constraint under a predeclared paired uncertainty test.  The dual variable starts
at $\Lambda=10^6$, making the initial path weight $10^{-6}$, and increases or decreases according
to a moving average of $c-\tau$.  A Boolean indicator that sparsity is enabled is not evidence
that pruning occurred.

The coefficient $\lambda_{\mathrm{logit}}$ is selected without mass labels.  The dense sweep
contains an exact zero control, rejects candidates whose prediction error is more than 5\% worse
than that control, and selects the smallest Pareto coefficient attaining 90\% of the best
admissible reduction in $\mathcal{L}_{\mathrm{logit}}-2$.  We repeat or revalidate this screen
whenever the token geometry, matching rule, or learned target space changes.  Neither masses nor
graph labels select hyperparameters, checkpoints, or seeds.

\paragraph{Parameter recovery.}
In Experiment~1, the learned parameter vector consists of five persistent global coordinates
$\widehat{\boldsymbol\theta}\in\mathbb{R}^{5}$.  On a probe-training fold we fit a scalar
one-hidden-layer MLP for every pair $(\widehat\theta_j,m_i)$.  A disjoint alignment fold yields
a nonlinear $R^2$ matrix and one global Hungarian bijection $\sigma$; the frozen bijection is
scored on a third fold:
\begin{equation}
    \operatorname{MCC}_{1:1}
      =\frac{1}{N}\sum_{i=1}^{N}R^2_{\sigma(i),i}.
    \label{eq:strict-mcc}
\end{equation}
Negative $R^2$ values are clipped to zero.  This is stricter than a mean--maximum score because
one learned coordinate cannot explain several masses and the permutation cannot change between
episodes.  The same frozen $\sigma$ aligns parameter columns before computing graph error.  In
the track-attached visual experiments, the single trajectory assignment instead pairs each
learned track scalar with its physical object; the nonlinear probe is then scored over held-out
episode--track pairs.

We report raw prediction error, the exact constraint in
Equation~\eqref{eq:raw-geco-constraint}, $\rho_{\mathrm{path}}$, state and parameter structural
Hamming distance (SHD), all five assigned nonlinear $R^2$ values, and
$\operatorname{MCC}_{1:1}$.  SHD is the mean unnormalised number of mismatched entries in the
$N\times N$ local graph and therefore lies in $[0,25]$.  Decodability alone does not establish
causal use.  We consequently evaluate three interventions: permuting learned parameter vectors
between episodes, replacing them by their training mean, and disabling retained
parameter-to-state paths.  Effects are reported separately on contact and non-contact
transitions.

Seed 0 is used for development.  Once a protocol is frozen, eight fresh paired seeds (1--8) are
reported, including valid optimization failures.  Dense, token-local, and sparse checkpoints are
evaluated on the same final 5,000-episode split.  Paired differences use a percentile bootstrap
over seed pairs with 10,000 resamples and a fixed analysis seed.  Interrupted jobs, corrupt
artifacts, or provenance mismatches are rerun; reproducible optimization failures remain in the
denominator.

\begin{table}[t]
    \centering
    \small
    \begin{tabular}{p{0.18\linewidth}p{0.22\linewidth}p{0.23\linewidth}p{0.27\linewidth}}
        \toprule
        Experiment & Predictor-side observation & Prediction target & Strongest supported claim \\
        \midrule
        1: true state & $Z_{0:t}$ & true $Z_{t+1}$ & Identification and sparse use conditional on ordered oracle state tracks. \\
        2: visual context & $X_{0:t}$ & true $Z_{t+1}$ & Visual state estimation and mass discovery with state-supervised training. \\
        3: fully visual & $X_{0:t}$ & $E_\phi(X_{t+1})$ & Jointly learned visual mass representation and sparse causal use, up to trajectory permutation. \\
        \bottomrule
    \end{tabular}
    \caption{The three experiments change the observation contract one step at a time.}
    \label{tab:experiment-ladder}
\end{table}

\subsection{Experiment 1: true-state context and true-state target}
\label{sec:experiments:true-state}

\paragraph{Question and source of mass information.}
This experiment asks whether the parameter encoder and SPARTAN can recover and use the masses
when perception, state grounding, and slot correspondence are given.  The model receives only
$[p_x,p_y,v_x,v_y]$ rows, but those trajectories contain mass information through collision
impulses, mass-dependent contact and wall geometry, and---in the current primary generator---the
mass-dependent initial-position support.  The mass-independent-initialization control isolates
the latter contribution.

\paragraph{Model.}
A shared linear map embeds each context state into 32 dimensions,
\begin{equation}
    e_t^i=W_c z_t^i+b_c\in\mathbb{R}^{32}.
\end{equation}
Five learned latent queries attend to all $T_{\mathrm{par}}N=150$ context tokens after learned
temporal and source-track embeddings, and a scalar head produces
\begin{equation}
    \widehat{S}^{\mathrm{ph}}
    \equiv\widehat{\boldsymbol\theta}
    =P_\eta(e_{0:T_{\mathrm{par}}-1})\in\mathbb{R}^{5\times1}.
    \label{eq:experiment-one-parameters}
\end{equation}
These are persistent global coordinates; coordinate $j$ is not declared to be the mass of input
row $j$.  Independent learned addresses distinguish the five state nodes and five parameter
nodes without privileging any of the 25 parameter-to-state relations.  SPARTAN separately
projects four-dimensional state tokens and scalar parameter tokens into a 512-dimensional
workspace, applies three sparse-attention layers with three-linear-layer 512-wide MLPs, and
projects the five decoded state tokens back to $\mathbb{R}^{4}$:
\begin{equation}
    \widehat Z_{t+1}
      =f_\gamma(Z_t,\widehat{\boldsymbol\theta}),
    \qquad t\in\mathcal I.
    \label{eq:experiment-one-predictor}
\end{equation}
The same $\widehat{\boldsymbol\theta}$ is reused for all 30 transitions in an episode.

Training uses aligned raw-state MSE,
\begin{equation}
    \mathcal{L}^{(1)}_{\mathrm{pred}}
    =\frac{1}{|\mathcal I|N\,4}
      \sum_{t\in\mathcal I}\sum_{i=1}^{N}
      \lVert\widehat z_{t+1}^i-z_{t+1}^i\rVert_2^2,
    \label{eq:experiment-one-loss}
\end{equation}
with every prediction anchored at the corresponding true $Z_t$.  The operative run uses 100,000
training episodes, batch size 16 (480 supervised transitions per update), Adam with learning rate
$5\times10^{-5}$, and 300,000 updates.  No target encoder is present.  VISReg is applied only to
learned context-side representations in the confirmatory protocol.

\paragraph{Current status.}
The state model, global parameter coordinates, dense and token-local references, learned hard
gates, raw constraint, one-to-one mass recovery, and graph evaluation are implemented.  Before a
paper claim, all three checkpoints must be evaluated on the same final split; evaluation artifacts
must be split-specific; paired bootstrap reporting and parameter interventions must be added; and
the gradient-free target VISReg constant must be removed from this regime.

\paragraph{Continuation gate.}
We begin implementation of Experiment~2 only after a complete development pilot shows that
(i) the dense constraint is reliably below the token-local constraint; (ii) the sparse constraint
crosses its freshly calibrated $\tau$; (iii) $\Lambda$ leaves its ceiling downwards and held-out
$\rho_{\mathrm{path}}$ falls for at least three consecutive evaluations; (iv) final density is
strictly between the token-local and dense endpoints; (v) five distinct masses are recovered
under one frozen global permutation; (vi) aligned parameter SHD improves; and (vii) at least one
parameter intervention worsens contact-sensitive prediction.  A failed but understood pilot may
motivate a protocol revision; adding pixels cannot repair an unexplained failure of this oracle
state experiment.

The confirmatory Experiment~1 claim additionally uses the following preregistered targets:
\begin{itemize}
    \item at least six of eight sparse seeds show sustained pruning and finish with the median of
          their last three density evaluations in $(0.10,0.90]$;
    \item the ratio of mean sparse to paired-dense test MSE is at most $1.05$, and at least six of
          eight seeds have a last-three-validation median constraint no larger than $1.05\tau$;
    \item median sparse $\operatorname{MCC}_{1:1}$ is at least $0.80$, with at least six of eight
          seeds at or above $0.70$;
    \item the paired 95\% interval for sparse minus token-local MCC lies above zero, and mean
          sparse parameter SHD is below both matched references;
    \item a positive paired sparse-minus-dense MCC interval is the strongest result, but a
          structural result remains supported when that interval overlaps zero if the ratio of
          mean sparse to dense MCC is at least $0.95$ and all prediction, pruning, graph, and
          intervention gates pass.
\end{itemize}
The Experiment~1 claim is conditional on ordered oracle kinematic tracks.  It is not yet a visual
or anonymous-object result, and high MCC without an intervention effect is described as parameter
information rather than demonstrated causal use.

\subsection{Permutation-safe readiness checks}
\label{sec:experiments:readiness}

The fixed global coordinates and node addresses in Experiment~1 are appropriate for ordered
simulator rows but not for anonymous visual slots.  Before introducing pixels, we validate the
visual coordinate contract on consistently permuted true-state trajectories.  The visual models
emit one scalar parameter per inferred track.  The state token and parameter token of track $i$
receive the same ephemeral key $\kappa_i$; $\kappa_i$ carries no state or physical information
and permutes with the complete track.  This binding is architectural, whereas the parameter value
and the SPARTAN paths that use it remain learned.

For a trajectory-level raw-state assignment, define the standardized cost
\begin{equation}
    D_e(i,j)=\frac{1}{|\mathcal I|\,4}
      \sum_{t\in\mathcal I}\sum_{a=1}^{4}
      \left(
        \frac{\widehat z_{e,t+1,i,a}-z_{e,t+1,j,a}}{\sigma_a}
      \right)^2,
    \qquad
    \pi_e^*=\underset{\pi\in\mathfrak S_N}{\operatorname{argmin}}
      \sum_{i=1}^{N}D_e(i,\pi(i)).
    \label{eq:trajectory-hungarian-raw}
\end{equation}
The discrete assignment is detached, but gradients pass through the selected predictions.  One
$\pi_e^*$ is held fixed for the complete trajectory and is reused for state loss, parameter
recovery, and every graph axis; per-frame rematching is forbidden.  The readiness suite requires
invariance to a whole-trajectory permutation, a positive penalty for a mid-trajectory switch,
joint equivariance of state tokens, parameter tokens and keys, and correct graph-axis alignment.
We retain exactly five slots until an explicit null-slot contract is introduced.  These are
implementation-validity checks, not a fourth scientific experiment.

\subsection{Experiment 2: visual context and true-state target}
\label{sec:experiments:visual-state}

\paragraph{Question and source of mass information.}
Experiment~2 asks whether a visual encoder can replace every predictor-side true state while the
prediction target remains the fixed physical state.  Because rendered glyphs have equal size and
carry no persistent simulator-row signature, mass must be inferred from a 30-frame history of
motion and contact geometry rather than from a single object's appearance.  The unchanged raw
target makes failures attributable to perception or parameter inference rather than to a moving
target representation.

\paragraph{Model.}
A causal recurrent encoder $Q_\psi$ maps frames to five 32-dimensional tracks,
\begin{equation}
    \widetilde S_{0:t}=Q_\psi(X_{0:t}),
    \qquad
    \widetilde S_t\in\mathbb{R}^{N\times32}.
\end{equation}
A shared grounded head gives the current physical-state estimate, while track-aware temporal and
cross-object pooling gives one scalar parameter per track:
\begin{align}
    Z_t^{\mathrm{ctx}} &=g_\omega(\widetilde S_t)\in\mathbb{R}^{N\times4},
    \\
    \widehat{S}^{\mathrm{ph}}
       &=P_\eta(\widetilde S_{0:T_{\mathrm{par}}-1})
         \in\mathbb{R}^{N\times1}.
    \label{eq:experiment-two-heads}
\end{align}
Only frames $0,\ldots,29$ may affect $\widehat S^{\mathrm{ph}}$.  Later source frames are
available for teacher-forcing $Z_t^{\mathrm{ctx}}$, but perturbing them must leave the parameter
estimate unchanged.  With the shared ephemeral track keys from
Section~\ref{sec:experiments:readiness}, SPARTAN uses a four-dimensional state interface, a scalar
parameter interface, and a 512-dimensional internal workspace:
\begin{equation}
    \widehat Z_{t+1}
      =f_\gamma(Z_t^{\mathrm{ctx}},\widehat S^{\mathrm{ph}})
      \in\mathbb{R}^{N\times4}.
    \label{eq:experiment-two-prediction}
\end{equation}
There is no target encoder in this experiment.  The target is the unchanged $Z_{t+1}$, and
training remains one-step teacher-forced; longer open-loop rollouts are evaluation diagnostics.

The trajectory assignment in Equation~\eqref{eq:trajectory-hungarian-raw} is selected once per
episode and reused for both current-state grounding and next-state prediction.  With targets
gathered under the detached assignment, the objective is
\begin{equation}
    \mathcal{L}^{(2)}
      =\mathcal{L}^{(2)}_{\mathrm{pred}}
       +\lambda_{\mathrm{ground}}\mathcal{L}_{\mathrm{ground}}
       +\lambda_{\mathrm{logit}}\mathcal{L}_{\mathrm{logit}}
       +\lambda_{\mathrm{reg}}\mathcal{R}(\widetilde S)
       +\Lambda^{-1}\mathcal{L}_{\mathrm{path}}.
    \label{eq:experiment-two-objective}
\end{equation}
The GECO constraint contains only raw next-state MSE and the weighted logit penalty, as in
Equation~\eqref{eq:raw-geco-constraint}; grounding and VISReg remain outside it.  There is no
target-side VISReg because the target is fixed data.

\paragraph{Continuation gate.}
Experiment~2 must satisfy all shared dense, token-local, sparse, pruning, graph, and intervention
checks under a freshly selected $\lambda_{\mathrm{logit}}$ and freshly calibrated $\tau$.  In
addition, we require: near-complete five-object coverage and a low trajectory switch rate,
including at contacts; accurate grounded position and velocity without target-frame leakage;
stronger held-out mass recovery from $\widehat S^{\mathrm{ph}}$ than incremental mass information
in the four-dimensional state bottleneck; no change in $\widehat S^{\mathrm{ph}}$ when frames after
$T_{\mathrm{par}}$ are perturbed; and success with mass-dependent physical radii but
mass-independent uniform rendering.  Tracking thresholds and grounding tolerances are frozen on
the development split before paired confirmatory seeds.

Only after a complete pilot meets this conjunction do we proceed to the fully visual target.
The supported claim is that frame inputs suffice at inference for state estimation and sparse mass
discovery under state-supervised training.  It is not a self-supervised result: true states define
the grounding loss, prediction target, and training assignment.  Removing the current-state
grounding is retained as an ablation, not as a separate main experiment.

\subsection{Optional diagnostic: learned target of the true future state}
\label{sec:experiments:state-target-diagnostic}

If Experiment~3 fails, we insert a small diagnostic between Experiments~2 and~3.  A jointly
learned shared per-object encoder maps the true $Z_{t+1}\in\mathbb{R}^{N\times4}$ to a
32-dimensional target, while the source remains visual.  Dense and token-local runs test whether
a simultaneously moving target can remain non-collapsed and comparable when object rows are
already given.  If this diagnostic succeeds but the image target fails, slot formation or target
tracking is implicated; if both fail, the moving target geometry or calibration is implicated.
Because this diagnostic still trains on true states and does not test single-frame visual slot
formation, it is not a mandatory experiment and does not automatically receive a full sparse
multi-seed study.

\subsection{Experiment 3: jointly learned visual context and target}
\label{sec:experiments:fully-visual}

\paragraph{Question.}
The final experiment asks whether mass recovery, parameter use, and SPARTAN pruning survive when
all trainable components are learned jointly from frames.  During training, true states, masses,
contacts, renderer masks, and simulator identities are unavailable to the model, loss, and
matching procedure.  They are used only for held-out evaluation.  The long-context encoder, the
single-frame observation/target encoder, parameter encoder, temporal lift, and SPARTAN are all
optimized simultaneously; an EMA or frozen target is a labelled control rather than the primary
model.

\paragraph{Independent single-frame observations and generic target tracks.}
The observation encoder is called independently on one image at a time,
\begin{equation}
    \{(o_t^r,M_t^r)\}_{r=1}^{N}=E_\phi(X_t),
    \qquad
    o_t^r\in\mathbb{R}^{32},
    \label{eq:single-frame-encoder}
\end{equation}
where $M_t^r$ is the slot assignment map.  In particular, the target call for transition
$t\rightarrow t+1$ receives only $X_{t+1}$.  The loss may associate outputs from several such
independent calls, but the encoder itself never receives a target sequence or context history.

Let $\mu(M)$ denote the spatial centroid of an assignment map.  Since the currently normalized
map sums to one for every slot, its total mass is not a valid presence score; any future
confidence statistic must instead be derived from pre-normalization assignment mass, logits, or
entropy.  Target tracks are assembled sequentially with a detached Hungarian assignment
\begin{equation}
    \rho_{e,t}^*
      =\underset{\rho\in\mathfrak S_N}{\operatorname{argmin}}
        \sum_{i=1}^{N}\mathcal{C}_{e,t}(i,\rho(i)),
    \qquad
    \overline o_{e,t}^{\,i}=o_{e,t}^{\rho_{e,t}^*(i)}.
    \label{eq:target-track-assignment}
\end{equation}
For the primary Bounce experiment, $\mathcal C$ is geometry-only: it combines the distance
between the candidate centroid and the constant-velocity prediction
$2\mu(M_{t-1}^i)-\mu(M_{t-2}^i)$ with overlap between the candidate support and a correspondingly
translated previous support.  The first two frames use the analogous position-only costs.  A
normalized learned-feature distance is a possible extension for more complex scenes, but it is
not a fixed renderer-identity key and must generalize under held-out appearance changes.  The
primary Bounce renderer uses uniform object glyphs, so no hand-specified appearance category can
solve Equation~\eqref{eq:target-track-assignment}.

\paragraph{A type-closed dynamical state.}
A single frame supplies position-like evidence but cannot determine velocity.  We therefore use a
shared two-frame lift,
\begin{equation}
    S_t^i=L_\xi\!\left([\nu(\overline o_{t-1}^{\,i}),
                         \nu(\overline o_t^{\,i})]\right)
       \in\mathbb{R}^{32},
    \label{eq:temporal-lift}
\end{equation}
where $\nu$ is a predeclared per-slot normalization.  In parallel, a separate causal recurrent
context encoder and track-aware parameter pooler use only the first 30 frames,
\begin{equation}
    \widetilde S_{0:T_{\mathrm{par}}-1}
       =Q_\psi(X_{0:T_{\mathrm{par}}-1}),
    \qquad
    \widehat S^{\mathrm{ph}}
       =P_\eta(\widetilde S_{0:T_{\mathrm{par}}-1})
       \in\mathbb{R}^{N\times1}.
    \label{eq:visual-parameter-context}
\end{equation}
The parameter and dynamical-state tokens use the same permutation-carrying track keys described
in Section~\ref{sec:experiments:readiness}.  SPARTAN works internally at width 512 and predicts
the next 32-dimensional observation slots,
\begin{equation}
    \widehat O_{t+1}=f_\gamma(S_t,\widehat S^{\mathrm{ph}})
       \in\mathbb{R}^{N\times32}.
    \label{eq:fully-visual-prediction}
\end{equation}
An autoregressive prediction is fed back through the same lift,
\begin{equation}
    \widehat S_{t+1}^{\,i}
       =L_\xi\!\left([\nu(\overline o_t^{\,i}),
                       \nu(\widehat o_{t+1}^{\,i})]\right),
    \qquad
    \widehat O_{t+2}=f_\gamma(\widehat S_{t+1},\widehat S^{\mathrm{ph}}).
    \label{eq:type-closed-rollout}
\end{equation}
This avoids feeding a target-space prediction directly into a context-head state whose equal
width does not guarantee equal meaning.

After independently encoded targets have been assembled into tracks, a second detached
trajectory-level Hungarian assignment $\pi_e^*$ matches predicted tracks to target tracks by
their mean normalized latent distance across all prediction times.  This one assignment is held
fixed for the complete loss and graph alignment.  The within-target assignments $\rho_{e,t}^*$
and the predictor-to-target assignment $\pi_e^*$ are distinct; replacing either with independent
per-frame prediction matching could hide track switches.

\paragraph{Objective and common-start comparison.}
Let $\mathcal{L}_{\mathrm{TF}}$ denote trajectory-matched one-step latent MSE and let
$\mathcal{L}^{(2)}_{\mathrm{AR}}$ denote the two-step loss obtained from
Equation~\eqref{eq:type-closed-rollout}.  The predictive objective is
\begin{equation}
    \mathcal{L}^{(3)}_{\mathrm{pred}}
      =\mathcal{L}_{\mathrm{TF}}
       +\beta\mathcal{L}^{(2)}_{\mathrm{AR}},
    \label{eq:fully-visual-predictive-loss}
\end{equation}
where $\beta$ is fixed before the reference runs.  Longer open-loop horizons through multiple
contacts are evaluated but are not initially backpropagated end to end.  VISReg is computed per
aligned track across episodes and time for both learned branches; flattening all slot addresses
together would not exclude a static slot-codebook collapse.

For the learned target, the dual uses a detached within-track variance scale,
\begin{align}
    V_{\mathrm{within}}
      &=\frac{1}{N\,32}\sum_{i=1}^{N}\sum_{a=1}^{32}
        \operatorname{Var}_{e,t}(\overline o_{e,t,i,a}),
        \\
    c_3
      &=\frac{\mathcal{L}^{(3)}_{\mathrm{pred}}}
              {\operatorname{sg}(V_{\mathrm{within}})}
        +\lambda_{\mathrm{logit}}\mathcal{L}_{\mathrm{logit}}
      \leq\tau_3.
    \label{eq:fully-visual-constraint}
\end{align}
This scale normalization does not by itself make two independently learned target geometries
comparable.  We therefore first train all representation modules with a dense SPARTAN from
scratch, then clone one common checkpoint into continued dense, token-local, and sparse legs.
The target and context encoders remain trainable in the primary sparse leg, and $\tau_3$ is
measured in the shared checkpoint geometry.  A frozen-target pruning leg is reported as a
fixed-ruler audit, not as a replacement for the fully joint result.  Throughout training we
measure target drift, within-track variance, effective rank, and fixed held-out probes so that a
smaller loss cannot be misread as equal prediction quality after target collapse or semantic
drift.

\paragraph{Final acceptance gate.}
The fully visual result is accepted only if all of the following hold:
\begin{itemize}
    \item every trainable branch receives gradients, and aligned context and target slots retain
          non-trivial within-track variance, effective rank, and object-position information;
    \item the target encoder is called on one frame at a time, later frames cannot affect
          $\widehat S^{\mathrm{ph}}$, and no true state, mass, contact, mask, or simulator identity
          enters the training assignment or loss;
    \item all five objects are covered and the target and predictor trajectory switch rates are
          negligible outside explicitly flagged observational ambiguities, including at contacts;
    \item position is decodable from one-frame slots, while velocity is substantially more
          decodable from the two-frame state in Equation~\eqref{eq:temporal-lift} than from one
          frame alone;
    \item the dense constraint is reliably below the token-local constraint in the common target
          geometry, and the sparse model reaches its fresh $\tau_3$, prunes materially, and
          remains strictly above the token-local density endpoint;
    \item masses are recoverable from $\widehat S^{\mathrm{ph}}$ but not from a fixed
          single-frame glyph signature, and parameter permutation, mean replacement, and retained
          path ablations worsen contact-sensitive prediction;
    \item open-loop latent error and norm are reported by horizon through multiple contacts, and
          dense, token-local, and sparse target representations remain sufficiently comparable for
          the constrained-loss comparison to be meaningful;
    \item the result reproduces over paired confirmatory seeds and under uniform and held-out
          heterogeneous-rendering controls.
\end{itemize}

For visually indistinguishable objects at a genuinely ambiguous crossing or contact, recovery of
an arbitrary simulator label is not identifiable from the observations.  In such cases we report
set-level mass recovery and ambiguity-aware tracking metrics rather than using a renderer-specific
identity shortcut.  Subject to this qualification, Experiment~3 supports the intended final
claim: jointly learned frame representations contain a track-attached, time-invariant mass
representation that SPARTAN uses through a pruned causal predictor, with true simulator variables
used only to evaluate that claim.
.
% Required packages: amsmath, amssymb, bm, booktabs.
% Status: Experiment 1 is implemented and running; Experiments 2 and 3 are
% prospective protocols. Replace the explicit status statements with measured
% results once each continuation gate has been passed.

\section{Experiments}
\label{sec:experiments}

We evaluate whether sparse predictive structure can identify time-invariant physical
parameters while progressively removing access to the ground-truth state.  We use a
three-experiment ladder.  Experiment~1 gives both branches the true physical state;
Experiment~2 replaces the predictor-side state by a causal visual encoding while retaining a
true-state target; and Experiment~3 jointly learns the context and single-frame target
representations from pixels.  The physical system, parameter distribution, SPARTAN predictor,
and evaluation logic are held fixed as far as the change in observation contract permits.
Each experiment is run only after its predecessor passes a predeclared continuation gate.  This
staging separates failures of causal parameter identification from failures of perception,
tracking, or jointly learned target geometry.

At the time of writing, Experiment~1 is implemented and its attention-logit sweep is running.
The later experiments below specify the intended confirmatory protocols and must not be read as
completed empirical results.

\subsection{Bounce}
\label{sec:experiments:bounce}

\subsubsection{Data}
\label{sec:experiments:bounce:data}

An episode contains $N=5$ elastic balls in the unit square.  We use $X_t$ for a rendered
frame, $Z_t=(z_t^1,\ldots,z_t^N)$ for the true kinematic state, and
$\boldsymbol{m}=(m_1,\ldots,m_N)$ for the time-invariant physical parameters, where
\begin{equation}
    X_t\in[0,1]^{3\times64\times64},
    \qquad
    z_t^i=(p_{x,t}^i,p_{y,t}^i,v_{x,t}^i,v_{y,t}^i)\in\mathbb{R}^{4},
    \qquad
    m_i>0.
    \label{eq:bounce-observations}
\end{equation}
For every episode and object, the implemented Baumgartner-aligned configuration samples
\begin{equation}
    \epsilon_i\sim\mathcal{N}(0,1),
    \qquad
    m_i=\operatorname{clip}_{[0.5,3.0]}(1.5+0.5\epsilon_i),
    \qquad
    r_i=r_0\frac{m_i}{m_{\mathrm{ref}}},
    \label{eq:bounce-mass-radius}
\end{equation}
with $r_0=0.08$ and the episode-independent reference $m_{\mathrm{ref}}=1.5$.
Equation~\eqref{eq:bounce-mass-radius} is a clipped Gaussian, rather than a truncated
Gaussian, and therefore places atoms at the clipping boundaries.  The numerical mass
distribution is our declared choice because its exact parameters are not specified by the
source Bounce experiment.

The initial velocity has fixed norm $0.7$ and uniformly random direction,
\begin{equation}
    \varphi_i\sim\mathcal{U}[0,2\pi),
    \qquad
    v_0^i=0.7(\cos\varphi_i,\sin\varphi_i).
\end{equation}
The current generator sequentially samples an initial centre uniformly from
$[1.05r_i,1-1.05r_i]^2$ and accepts it only if
$\lVert p_0^i-p_0^j\rVert_2>1.1(r_i+r_j)$ for all previously placed objects $j$.
Consequently, the support of the initial positions is itself mass-dependent.  We report this
fact rather than describing the initial state as independent of mass.  To isolate evidence from
the subsequent dynamics, the final study also includes a mass-independent-initialization
control: all wall margins and pair-separation constraints are computed using the common
worst-case radius $r_{\max}=r_0(3.0/1.5)$, irrespective of the sampled masses.  This control is
required before attributing recovery specifically to transition dynamics rather than to the
initial-state support.

The simulator records $T=60$ states at intervals $\Delta t=0.1$, with eight physics
substeps per recorded transition.  Free motion is followed by exact reflection of wall
overshoot about $r_i$ or $1-r_i$.  For a colliding pair, let
\begin{equation}
    n_{ij}=\frac{p_i-p_j}{\lVert p_i-p_j\rVert_2},
    \qquad
    u_{ij}=(v_i-v_j)^\top n_{ij}<0.
\end{equation}
The restitution-one velocity update is
\begin{align}
    v_i' &= v_i-\frac{2m_j}{m_i+m_j}u_{ij}n_{ij},
    &
    v_j' &= v_j+\frac{2m_i}{m_i+m_j}u_{ij}n_{ij}.
    \label{eq:elastic-collision}
\end{align}
Before applying this impulse, any penetration is projected back to the contact surface
$\lVert p_i-p_j\rVert_2=r_i+r_j$ using inverse-mass weights.  The corrected simulator is
versioned so that pre-generated trajectories from older collision rules cannot be mixed with
the current data.

Each episode supplies
\begin{equation}
    (X_{0:T-1},Z_{0:T-1},\boldsymbol{m},C_{0:T-2}),
    \qquad
    C_t\in\{0,1\}^{N\times N},
\end{equation}
where $C_{t,ij}=1$ records a pair contact during transition $t\rightarrow t+1$.
Off-diagonal entries are symmetric.  In the mass-dependent-radius environment,
$C_{t,ii}=1$ records a wall contact because its location depends on $m_i$ through $r_i$.
We use the first $T_{\mathrm{par}}=30$ observations to infer one parameter representation and
evaluate $K=T-T_{\mathrm{par}}=30$ transitions,
\begin{equation}
    \mathcal{I}=\{T_{\mathrm{par}}-1,\ldots,T-2\}=\{29,\ldots,58\}.
    \label{eq:transition-index-set}
\end{equation}
Thus, the parameter encoder sees observations $0,\ldots,29$, while the supervised pairs in
Experiment~1 are $(Z_{29},Z_{30}),\ldots,(Z_{58},Z_{59})$.  These are 30
teacher-forced one-step predictions, not one 30-step open-loop rollout.

\paragraph{Why absolute masses are observable.}
If all physical radii were equal, multiplying every mass by a common positive constant would
leave Equation~\eqref{eq:elastic-collision} unchanged; elastic impulses would then identify
mass ratios but not absolute scale.  The implemented relation
$r_i=r_0m_i/m_{\mathrm{ref}}$ removes that symmetry because a common mass rescaling also
changes pair-contact distances, wall-contact locations, penetration correction, and the support
of initial positions.  The current state experiment sets rendering off and never supplies a
literal radius coordinate, but its observed positions and velocities are nevertheless generated
by this mass-dependent geometry.  We therefore claim identification in the
mass-proportional-physical-radius environment, not in the harder equal-radius environment.

For Experiments~2 and~3, physical and rendered radii must be separated.  The simulator retains
Equation~\eqref{eq:bounce-mass-radius}, while every object is drawn with the same
mass-independent glyph and rendered radius.  The primary visual data must also contain no
dataset-global visual signature tied to simulator row index.  Hence a single frame cannot reveal
mass from the object's glyph, although a sequence can reveal mass through centre motion and
hidden contact geometry.  A visible-radius version and an equal-physical-radius version are
reported only as labelled controls.  The current data implementation uses a single radius switch
for physics and rendering; splitting those two settings and versioning the resulting data are
prerequisites for the visual experiments.

\subsubsection{Ground-truth local graphs}
\label{sec:experiments:bounce:graphs}

The destination row is indexed first in every graph.  From $C_t$, we define the true
state-to-state and mass-to-state local graphs as
\begin{align}
    G^{Z\rightarrow Z}_{t,ij}
        &= \mathbf{1}\{i=j\ \lor\ C_{t,ij}=1\},
        \\
    G^{m\rightarrow Z}_{t,ij}
        &= \mathbf{1}\!\left\{C_{t,ij}=1\ \lor\
             \left(i=j\ \land\ \bigvee_k C_{t,ik}=1\right)\right\}.
    \label{eq:bounce-ground-truth-graphs}
\end{align}
A pair collision therefore gives each affected state an own-mass and partner-mass parent; a
mass-dependent wall collision gives the affected state an own-mass parent.  Self-state edges are
always present because free motion depends on the object's previous state.

\subsubsection{Shared model comparisons and metrics}
\label{sec:experiments:shared-protocol}

Every scientific comparison uses the same three predictor modes under matched initialization
rules, training budgets, and data splits:
\begin{enumerate}
    \item \emph{Dense}: $A^{(\ell)}\equiv\boldsymbol{1}$ in every SPARTAN layer and no
          path penalty.
    \item \emph{Token-local}: $A^{(\ell)}\equiv\boldsymbol{0}$, so only each token's
          residual--MLP path remains.  This reference has no access to other state or parameter
          tokens, but its own state may contain statistical traces of earlier mass-dependent
          geometry; it is therefore more precise than calling the reference fully mass-blind.
    \item \emph{Sparse}: learned hard Bernoulli gates with the SPARTAN path objective and a
          GECO-style constraint controller.
\end{enumerate}

There are $N$ state tokens and $N$ parameter tokens.  For $L_{\mathrm{sp}}=3$ SPARTAN
layers, the multilayer reachability matrix is
\begin{equation}
    \overline A
    =\bigl(A^{(L_{\mathrm{sp}})}+I\bigr)\cdots
     \bigl(A^{(1)}+I\bigr),
    \label{eq:spartan-path-matrix}
\end{equation}
where $\overline A_{ij}>0$ means that at least one permitted computational path connects source
token $j$ to decoded output $i$.  The optimized path count includes only the $N$ decoded state
rows,
\begin{equation}
    \mathcal{L}_{\mathrm{path}}
    =\mathbb{E}\!\left[
       \sum_{i=1}^{N}\sum_{j=1}^{2N}\overline A_{ij}
     \right],
    \qquad
    \rho_{\mathrm{path}}
    =\frac{1}{2N^2}\sum_{i=1}^{N}\sum_{j=1}^{2N}
       \mathbf{1}\{\overline A_{ij}>0\}.
    \label{eq:path-objective-density}
\end{equation}
The second quantity discards path multiplicity and is the primary pruning metric.  With ten
tokens, the token-local reference has $\overline A=I_{10}$ and
$\rho_{\mathrm{path}}=5/50=0.1$, while the dense model has
$\rho_{\mathrm{path}}=1$.  The corresponding unnormalised path objectives are 5 and 6655:
if $J$ is the $10\times10$ all-ones matrix, then
$(J+I)^3=I+133J$, each decoded row sums to 1331, and five decoded rows sum to 6655.
Reachability over all output rows is retained as a diagnostic only, since parameter-token outputs
are not decoded.

For a fixed target space, training minimizes
\begin{equation}
    \mathcal{L}
    =\mathcal{L}_{\mathrm{pred}}
     +\lambda_{\mathrm{logit}}\mathcal{L}_{\mathrm{logit}}
     +\lambda_{\mathrm{reg}}\mathcal{R}
     +\Lambda^{-1}\mathcal{L}_{\mathrm{path}},
    \label{eq:shared-training-objective}
\end{equation}
when sparsity is enabled.  Here $\mathcal{R}$ is VISReg on learned
representations.  Fixed raw-state targets are not regularized: applying VISReg to them produces
no model gradient and only adds a stochastic constant.  The attention-logit penalty averages
$\exp(a)+\exp(-a)$ over layer logits $a$ and has theoretical minimum 2; the implementation uses
a linear continuation beyond $|a|=30$ for numerical stability.

The dual constraint excludes both VISReg and the path penalty,
\begin{equation}
    c=\mathcal{L}_{\mathrm{pred}}
      +\lambda_{\mathrm{logit}}\mathcal{L}_{\mathrm{logit}}
      \leq \tau.
    \label{eq:raw-geco-constraint}
\end{equation}
For every architecture and seed, $\tau$ is recalibrated as the held-out constraint of a
converged dense reference, and the sparse run is launched only if the dense constraint is below
the token-local constraint under a predeclared paired uncertainty test.  The dual variable starts
at $\Lambda=10^6$, making the initial path weight $10^{-6}$, and increases or decreases according
to a moving average of $c-\tau$.  A Boolean indicator that sparsity is enabled is not evidence
that pruning occurred.

The coefficient $\lambda_{\mathrm{logit}}$ is selected without mass labels.  The dense sweep
contains an exact zero control, rejects candidates whose prediction error is more than 5\% worse
than that control, and selects the smallest Pareto coefficient attaining 90\% of the best
admissible reduction in $\mathcal{L}_{\mathrm{logit}}-2$.  We repeat or revalidate this screen
whenever the token geometry, matching rule, or learned target space changes.  Neither masses nor
graph labels select hyperparameters, checkpoints, or seeds.

\paragraph{Parameter recovery.}
In Experiment~1, the learned parameter vector consists of five persistent global coordinates
$\widehat{\boldsymbol\theta}\in\mathbb{R}^{5}$.  On a probe-training fold we fit a scalar
one-hidden-layer MLP for every pair $(\widehat\theta_j,m_i)$.  A disjoint alignment fold yields
a nonlinear $R^2$ matrix and one global Hungarian bijection $\sigma$; the frozen bijection is
scored on a third fold:
\begin{equation}
    \operatorname{MCC}_{1:1}
      =\frac{1}{N}\sum_{i=1}^{N}R^2_{\sigma(i),i}.
    \label{eq:strict-mcc}
\end{equation}
Negative $R^2$ values are clipped to zero.  This is stricter than a mean--maximum score because
one learned coordinate cannot explain several masses and the permutation cannot change between
episodes.  The same frozen $\sigma$ aligns parameter columns before computing graph error.  In
the track-attached visual experiments, the single trajectory assignment instead pairs each
learned track scalar with its physical object; the nonlinear probe is then scored over held-out
episode--track pairs.

We report raw prediction error, the exact constraint in
Equation~\eqref{eq:raw-geco-constraint}, $\rho_{\mathrm{path}}$, state and parameter structural
Hamming distance (SHD), all five assigned nonlinear $R^2$ values, and
$\operatorname{MCC}_{1:1}$.  SHD is the mean unnormalised number of mismatched entries in the
$N\times N$ local graph and therefore lies in $[0,25]$.  Decodability alone does not establish
causal use.  We consequently evaluate three interventions: permuting learned parameter vectors
between episodes, replacing them by their training mean, and disabling retained
parameter-to-state paths.  Effects are reported separately on contact and non-contact
transitions.

Seed 0 is used for development.  Once a protocol is frozen, eight fresh paired seeds (1--8) are
reported, including valid optimization failures.  Dense, token-local, and sparse checkpoints are
evaluated on the same final 5,000-episode split.  Paired differences use a percentile bootstrap
over seed pairs with 10,000 resamples and a fixed analysis seed.  Interrupted jobs, corrupt
artifacts, or provenance mismatches are rerun; reproducible optimization failures remain in the
denominator.

\begin{table}[t]
    \centering
    \small
    \begin{tabular}{p{0.18\linewidth}p{0.22\linewidth}p{0.23\linewidth}p{0.27\linewidth}}
        \toprule
        Experiment & Predictor-side observation & Prediction target & Strongest supported claim \\
        \midrule
        1: true state & $Z_{0:t}$ & true $Z_{t+1}$ & Identification and sparse use conditional on ordered oracle state tracks. \\
        2: visual context & $X_{0:t}$ & true $Z_{t+1}$ & Visual state estimation and mass discovery with state-supervised training. \\
        3: fully visual & $X_{0:t}$ & $E_\phi(X_{t+1})$ & Jointly learned visual mass representation and sparse causal use, up to trajectory permutation. \\
        \bottomrule
    \end{tabular}
    \caption{The three experiments change the observation contract one step at a time.}
    \label{tab:experiment-ladder}
\end{table}

\subsection{Experiment 1: true-state context and true-state target}
\label{sec:experiments:true-state}

\paragraph{Question and source of mass information.}
This experiment asks whether the parameter encoder and SPARTAN can recover and use the masses
when perception, state grounding, and slot correspondence are given.  The model receives only
$[p_x,p_y,v_x,v_y]$ rows, but those trajectories contain mass information through collision
impulses, mass-dependent contact and wall geometry, and---in the current primary generator---the
mass-dependent initial-position support.  The mass-independent-initialization control isolates
the latter contribution.

\paragraph{Model.}
A shared linear map embeds each context state into 32 dimensions,
\begin{equation}
    e_t^i=W_c z_t^i+b_c\in\mathbb{R}^{32}.
\end{equation}
Five learned latent queries attend to all $T_{\mathrm{par}}N=150$ context tokens after learned
temporal and source-track embeddings, and a scalar head produces
\begin{equation}
    \widehat{S}^{\mathrm{ph}}
    \equiv\widehat{\boldsymbol\theta}
    =P_\eta(e_{0:T_{\mathrm{par}}-1})\in\mathbb{R}^{5\times1}.
    \label{eq:experiment-one-parameters}
\end{equation}
These are persistent global coordinates; coordinate $j$ is not declared to be the mass of input
row $j$.  Independent learned addresses distinguish the five state nodes and five parameter
nodes without privileging any of the 25 parameter-to-state relations.  SPARTAN separately
projects four-dimensional state tokens and scalar parameter tokens into a 512-dimensional
workspace, applies three sparse-attention layers with three-linear-layer 512-wide MLPs, and
projects the five decoded state tokens back to $\mathbb{R}^{4}$:
\begin{equation}
    \widehat Z_{t+1}
      =f_\gamma(Z_t,\widehat{\boldsymbol\theta}),
    \qquad t\in\mathcal I.
    \label{eq:experiment-one-predictor}
\end{equation}
The same $\widehat{\boldsymbol\theta}$ is reused for all 30 transitions in an episode.

Training uses aligned raw-state MSE,
\begin{equation}
    \mathcal{L}^{(1)}_{\mathrm{pred}}
    =\frac{1}{|\mathcal I|N\,4}
      \sum_{t\in\mathcal I}\sum_{i=1}^{N}
      \lVert\widehat z_{t+1}^i-z_{t+1}^i\rVert_2^2,
    \label{eq:experiment-one-loss}
\end{equation}
with every prediction anchored at the corresponding true $Z_t$.  The operative run uses 100,000
training episodes, batch size 16 (480 supervised transitions per update), Adam with learning rate
$5\times10^{-5}$, and 300,000 updates.  No target encoder is present.  VISReg is applied only to
learned context-side representations in the confirmatory protocol.

\paragraph{Current status.}
The state model, global parameter coordinates, dense and token-local references, learned hard
gates, raw constraint, one-to-one mass recovery, and graph evaluation are implemented.  Before a
paper claim, all three checkpoints must be evaluated on the same final split; evaluation artifacts
must be split-specific; paired bootstrap reporting and parameter interventions must be added; and
the gradient-free target VISReg constant must be removed from this regime.

\paragraph{Continuation gate.}
We begin implementation of Experiment~2 only after a complete development pilot shows that
(i) the dense constraint is reliably below the token-local constraint; (ii) the sparse constraint
crosses its freshly calibrated $\tau$; (iii) $\Lambda$ leaves its ceiling downwards and held-out
$\rho_{\mathrm{path}}$ falls for at least three consecutive evaluations; (iv) final density is
strictly between the token-local and dense endpoints; (v) five distinct masses are recovered
under one frozen global permutation; (vi) aligned parameter SHD improves; and (vii) at least one
parameter intervention worsens contact-sensitive prediction.  A failed but understood pilot may
motivate a protocol revision; adding pixels cannot repair an unexplained failure of this oracle
state experiment.

The confirmatory Experiment~1 claim additionally uses the following preregistered targets:
\begin{itemize}
    \item at least six of eight sparse seeds show sustained pruning and finish with the median of
          their last three density evaluations in $(0.10,0.90]$;
    \item the ratio of mean sparse to paired-dense test MSE is at most $1.05$, and at least six of
          eight seeds have a last-three-validation median constraint no larger than $1.05\tau$;
    \item median sparse $\operatorname{MCC}_{1:1}$ is at least $0.80$, with at least six of eight
          seeds at or above $0.70$;
    \item the paired 95\% interval for sparse minus token-local MCC lies above zero, and mean
          sparse parameter SHD is below both matched references;
    \item a positive paired sparse-minus-dense MCC interval is the strongest result, but a
          structural result remains supported when that interval overlaps zero if the ratio of
          mean sparse to dense MCC is at least $0.95$ and all prediction, pruning, graph, and
          intervention gates pass.
\end{itemize}
The Experiment~1 claim is conditional on ordered oracle kinematic tracks.  It is not yet a visual
or anonymous-object result, and high MCC without an intervention effect is described as parameter
information rather than demonstrated causal use.

\subsection{Permutation-safe readiness checks}
\label{sec:experiments:readiness}

The fixed global coordinates and node addresses in Experiment~1 are appropriate for ordered
simulator rows but not for anonymous visual slots.  Before introducing pixels, we validate the
visual coordinate contract on consistently permuted true-state trajectories.  The visual models
emit one scalar parameter per inferred track.  The state token and parameter token of track $i$
receive the same ephemeral key $\kappa_i$; $\kappa_i$ carries no state or physical information
and permutes with the complete track.  This binding is architectural, whereas the parameter value
and the SPARTAN paths that use it remain learned.

For a trajectory-level raw-state assignment, define the standardized cost
\begin{equation}
    D_e(i,j)=\frac{1}{|\mathcal I|\,4}
      \sum_{t\in\mathcal I}\sum_{a=1}^{4}
      \left(
        \frac{\widehat z_{e,t+1,i,a}-z_{e,t+1,j,a}}{\sigma_a}
      \right)^2,
    \qquad
    \pi_e^*=\underset{\pi\in\mathfrak S_N}{\operatorname{argmin}}
      \sum_{i=1}^{N}D_e(i,\pi(i)).
    \label{eq:trajectory-hungarian-raw}
\end{equation}
The discrete assignment is detached, but gradients pass through the selected predictions.  One
$\pi_e^*$ is held fixed for the complete trajectory and is reused for state loss, parameter
recovery, and every graph axis; per-frame rematching is forbidden.  The readiness suite requires
invariance to a whole-trajectory permutation, a positive penalty for a mid-trajectory switch,
joint equivariance of state tokens, parameter tokens and keys, and correct graph-axis alignment.
We retain exactly five slots until an explicit null-slot contract is introduced.  These are
implementation-validity checks, not a fourth scientific experiment.

\subsection{Experiment 2: visual context and true-state target}
\label{sec:experiments:visual-state}

\paragraph{Question and source of mass information.}
Experiment~2 asks whether a visual encoder can replace every predictor-side true state while the
prediction target remains the fixed physical state.  Because rendered glyphs have equal size and
carry no persistent simulator-row signature, mass must be inferred from a 30-frame history of
motion and contact geometry rather than from a single object's appearance.  The unchanged raw
target makes failures attributable to perception or parameter inference rather than to a moving
target representation.

\paragraph{Model.}
A causal recurrent encoder $Q_\psi$ maps frames to five 32-dimensional tracks,
\begin{equation}
    \widetilde S_{0:t}=Q_\psi(X_{0:t}),
    \qquad
    \widetilde S_t\in\mathbb{R}^{N\times32}.
\end{equation}
A shared grounded head gives the current physical-state estimate, while track-aware temporal and
cross-object pooling gives one scalar parameter per track:
\begin{align}
    Z_t^{\mathrm{ctx}} &=g_\omega(\widetilde S_t)\in\mathbb{R}^{N\times4},
    \\
    \widehat{S}^{\mathrm{ph}}
       &=P_\eta(\widetilde S_{0:T_{\mathrm{par}}-1})
         \in\mathbb{R}^{N\times1}.
    \label{eq:experiment-two-heads}
\end{align}
Only frames $0,\ldots,29$ may affect $\widehat S^{\mathrm{ph}}$.  Later source frames are
available for teacher-forcing $Z_t^{\mathrm{ctx}}$, but perturbing them must leave the parameter
estimate unchanged.  With the shared ephemeral track keys from
Section~\ref{sec:experiments:readiness}, SPARTAN uses a four-dimensional state interface, a scalar
parameter interface, and a 512-dimensional internal workspace:
\begin{equation}
    \widehat Z_{t+1}
      =f_\gamma(Z_t^{\mathrm{ctx}},\widehat S^{\mathrm{ph}})
      \in\mathbb{R}^{N\times4}.
    \label{eq:experiment-two-prediction}
\end{equation}
There is no target encoder in this experiment.  The target is the unchanged $Z_{t+1}$, and
training remains one-step teacher-forced; longer open-loop rollouts are evaluation diagnostics.

The trajectory assignment in Equation~\eqref{eq:trajectory-hungarian-raw} is selected once per
episode and reused for both current-state grounding and next-state prediction.  With targets
gathered under the detached assignment, the objective is
\begin{equation}
    \mathcal{L}^{(2)}
      =\mathcal{L}^{(2)}_{\mathrm{pred}}
       +\lambda_{\mathrm{ground}}\mathcal{L}_{\mathrm{ground}}
       +\lambda_{\mathrm{logit}}\mathcal{L}_{\mathrm{logit}}
       +\lambda_{\mathrm{reg}}\mathcal{R}(\widetilde S)
       +\Lambda^{-1}\mathcal{L}_{\mathrm{path}}.
    \label{eq:experiment-two-objective}
\end{equation}
The GECO constraint contains only raw next-state MSE and the weighted logit penalty, as in
Equation~\eqref{eq:raw-geco-constraint}; grounding and VISReg remain outside it.  There is no
target-side VISReg because the target is fixed data.

\paragraph{Continuation gate.}
Experiment~2 must satisfy all shared dense, token-local, sparse, pruning, graph, and intervention
checks under a freshly selected $\lambda_{\mathrm{logit}}$ and freshly calibrated $\tau$.  In
addition, we require: near-complete five-object coverage and a low trajectory switch rate,
including at contacts; accurate grounded position and velocity without target-frame leakage;
stronger held-out mass recovery from $\widehat S^{\mathrm{ph}}$ than incremental mass information
in the four-dimensional state bottleneck; no change in $\widehat S^{\mathrm{ph}}$ when frames after
$T_{\mathrm{par}}$ are perturbed; and success with mass-dependent physical radii but
mass-independent uniform rendering.  Tracking thresholds and grounding tolerances are frozen on
the development split before paired confirmatory seeds.

Only after a complete pilot meets this conjunction do we proceed to the fully visual target.
The supported claim is that frame inputs suffice at inference for state estimation and sparse mass
discovery under state-supervised training.  It is not a self-supervised result: true states define
the grounding loss, prediction target, and training assignment.  Removing the current-state
grounding is retained as an ablation, not as a separate main experiment.

\subsection{Optional diagnostic: learned target of the true future state}
\label{sec:experiments:state-target-diagnostic}

If Experiment~3 fails, we insert a small diagnostic between Experiments~2 and~3.  A jointly
learned shared per-object encoder maps the true $Z_{t+1}\in\mathbb{R}^{N\times4}$ to a
32-dimensional target, while the source remains visual.  Dense and token-local runs test whether
a simultaneously moving target can remain non-collapsed and comparable when object rows are
already given.  If this diagnostic succeeds but the image target fails, slot formation or target
tracking is implicated; if both fail, the moving target geometry or calibration is implicated.
Because this diagnostic still trains on true states and does not test single-frame visual slot
formation, it is not a mandatory experiment and does not automatically receive a full sparse
multi-seed study.

\subsection{Experiment 3: jointly learned visual context and target}
\label{sec:experiments:fully-visual}

\paragraph{Question.}
The final experiment asks whether mass recovery, parameter use, and SPARTAN pruning survive when
all trainable components are learned jointly from frames.  During training, true states, masses,
contacts, renderer masks, and simulator identities are unavailable to the model, loss, and
matching procedure.  They are used only for held-out evaluation.  The long-context encoder, the
single-frame observation/target encoder, parameter encoder, temporal lift, and SPARTAN are all
optimized simultaneously; an EMA or frozen target is a labelled control rather than the primary
model.

\paragraph{Independent single-frame observations and generic target tracks.}
The observation encoder is called independently on one image at a time,
\begin{equation}
    \{(o_t^r,M_t^r)\}_{r=1}^{N}=E_\phi(X_t),
    \qquad
    o_t^r\in\mathbb{R}^{32},
    \label{eq:single-frame-encoder}
\end{equation}
where $M_t^r$ is the slot assignment map.  In particular, the target call for transition
$t\rightarrow t+1$ receives only $X_{t+1}$.  The loss may associate outputs from several such
independent calls, but the encoder itself never receives a target sequence or context history.

Let $\mu(M)$ denote the spatial centroid of an assignment map.  Since the currently normalized
map sums to one for every slot, its total mass is not a valid presence score; any future
confidence statistic must instead be derived from pre-normalization assignment mass, logits, or
entropy.  Target tracks are assembled sequentially with a detached Hungarian assignment
\begin{equation}
    \rho_{e,t}^*
      =\underset{\rho\in\mathfrak S_N}{\operatorname{argmin}}
        \sum_{i=1}^{N}\mathcal{C}_{e,t}(i,\rho(i)),
    \qquad
    \overline o_{e,t}^{\,i}=o_{e,t}^{\rho_{e,t}^*(i)}.
    \label{eq:target-track-assignment}
\end{equation}
For the primary Bounce experiment, $\mathcal C$ is geometry-only: it combines the distance
between the candidate centroid and the constant-velocity prediction
$2\mu(M_{t-1}^i)-\mu(M_{t-2}^i)$ with overlap between the candidate support and a correspondingly
translated previous support.  The first two frames use the analogous position-only costs.  A
normalized learned-feature distance is a possible extension for more complex scenes, but it is
not a fixed renderer-identity key and must generalize under held-out appearance changes.  The
primary Bounce renderer uses uniform object glyphs, so no hand-specified appearance category can
solve Equation~\eqref{eq:target-track-assignment}.

\paragraph{A type-closed dynamical state.}
A single frame supplies position-like evidence but cannot determine velocity.  We therefore use a
shared two-frame lift,
\begin{equation}
    S_t^i=L_\xi\!\left([\nu(\overline o_{t-1}^{\,i}),
                         \nu(\overline o_t^{\,i})]\right)
       \in\mathbb{R}^{32},
    \label{eq:temporal-lift}
\end{equation}
where $\nu$ is a predeclared per-slot normalization.  In parallel, a separate causal recurrent
context encoder and track-aware parameter pooler use only the first 30 frames,
\begin{equation}
    \widetilde S_{0:T_{\mathrm{par}}-1}
       =Q_\psi(X_{0:T_{\mathrm{par}}-1}),
    \qquad
    \widehat S^{\mathrm{ph}}
       =P_\eta(\widetilde S_{0:T_{\mathrm{par}}-1})
       \in\mathbb{R}^{N\times1}.
    \label{eq:visual-parameter-context}
\end{equation}
The parameter and dynamical-state tokens use the same permutation-carrying track keys described
in Section~\ref{sec:experiments:readiness}.  SPARTAN works internally at width 512 and predicts
the next 32-dimensional observation slots,
\begin{equation}
    \widehat O_{t+1}=f_\gamma(S_t,\widehat S^{\mathrm{ph}})
       \in\mathbb{R}^{N\times32}.
    \label{eq:fully-visual-prediction}
\end{equation}
An autoregressive prediction is fed back through the same lift,
\begin{equation}
    \widehat S_{t+1}^{\,i}
       =L_\xi\!\left([\nu(\overline o_t^{\,i}),
                       \nu(\widehat o_{t+1}^{\,i})]\right),
    \qquad
    \widehat O_{t+2}=f_\gamma(\widehat S_{t+1},\widehat S^{\mathrm{ph}}).
    \label{eq:type-closed-rollout}
\end{equation}
This avoids feeding a target-space prediction directly into a context-head state whose equal
width does not guarantee equal meaning.

After independently encoded targets have been assembled into tracks, a second detached
trajectory-level Hungarian assignment $\pi_e^*$ matches predicted tracks to target tracks by
their mean normalized latent distance across all prediction times.  This one assignment is held
fixed for the complete loss and graph alignment.  The within-target assignments $\rho_{e,t}^*$
and the predictor-to-target assignment $\pi_e^*$ are distinct; replacing either with independent
per-frame prediction matching could hide track switches.

\paragraph{Objective and common-start comparison.}
Let $\mathcal{L}_{\mathrm{TF}}$ denote trajectory-matched one-step latent MSE and let
$\mathcal{L}^{(2)}_{\mathrm{AR}}$ denote the two-step loss obtained from
Equation~\eqref{eq:type-closed-rollout}.  The predictive objective is
\begin{equation}
    \mathcal{L}^{(3)}_{\mathrm{pred}}
      =\mathcal{L}_{\mathrm{TF}}
       +\beta\mathcal{L}^{(2)}_{\mathrm{AR}},
    \label{eq:fully-visual-predictive-loss}
\end{equation}
where $\beta$ is fixed before the reference runs.  Longer open-loop horizons through multiple
contacts are evaluated but are not initially backpropagated end to end.  VISReg is computed per
aligned track across episodes and time for both learned branches; flattening all slot addresses
together would not exclude a static slot-codebook collapse.

For the learned target, the dual uses a detached within-track variance scale,
\begin{align}
    V_{\mathrm{within}}
      &=\frac{1}{N\,32}\sum_{i=1}^{N}\sum_{a=1}^{32}
        \operatorname{Var}_{e,t}(\overline o_{e,t,i,a}),
        \\
    c_3
      &=\frac{\mathcal{L}^{(3)}_{\mathrm{pred}}}
              {\operatorname{sg}(V_{\mathrm{within}})}
        +\lambda_{\mathrm{logit}}\mathcal{L}_{\mathrm{logit}}
      \leq\tau_3.
    \label{eq:fully-visual-constraint}
\end{align}
This scale normalization does not by itself make two independently learned target geometries
comparable.  We therefore first train all representation modules with a dense SPARTAN from
scratch, then clone one common checkpoint into continued dense, token-local, and sparse legs.
The target and context encoders remain trainable in the primary sparse leg, and $\tau_3$ is
measured in the shared checkpoint geometry.  A frozen-target pruning leg is reported as a
fixed-ruler audit, not as a replacement for the fully joint result.  Throughout training we
measure target drift, within-track variance, effective rank, and fixed held-out probes so that a
smaller loss cannot be misread as equal prediction quality after target collapse or semantic
drift.

\paragraph{Final acceptance gate.}
The fully visual result is accepted only if all of the following hold:
\begin{itemize}
    \item every trainable branch receives gradients, and aligned context and target slots retain
          non-trivial within-track variance, effective rank, and object-position information;
    \item the target encoder is called on one frame at a time, later frames cannot affect
          $\widehat S^{\mathrm{ph}}$, and no true state, mass, contact, mask, or simulator identity
          enters the training assignment or loss;
    \item all five objects are covered and the target and predictor trajectory switch rates are
          negligible outside explicitly flagged observational ambiguities, including at contacts;
    \item position is decodable from one-frame slots, while velocity is substantially more
          decodable from the two-frame state in Equation~\eqref{eq:temporal-lift} than from one
          frame alone;
    \item the dense constraint is reliably below the token-local constraint in the common target
          geometry, and the sparse model reaches its fresh $\tau_3$, prunes materially, and
          remains strictly above the token-local density endpoint;
    \item masses are recoverable from $\widehat S^{\mathrm{ph}}$ but not from a fixed
          single-frame glyph signature, and parameter permutation, mean replacement, and retained
          path ablations worsen contact-sensitive prediction;
    \item open-loop latent error and norm are reported by horizon through multiple contacts, and
          dense, token-local, and sparse target representations remain sufficiently comparable for
          the constrained-loss comparison to be meaningful;
    \item the result reproduces over paired confirmatory seeds and under uniform and held-out
          heterogeneous-rendering controls.
\end{itemize}

For visually indistinguishable objects at a genuinely ambiguous crossing or contact, recovery of
an arbitrary simulator label is not identifiable from the observations.  In such cases we report
set-level mass recovery and ambiguity-aware tracking metrics rather than using a renderer-specific
identity shortcut.  Subject to this qualification, Experiment~3 supports the intended final
claim: jointly learned frame representations contain a track-attached, time-invariant mass
representation that SPARTAN uses through a pruned causal predictor, with true simulator variables
used only to evaluate that claim.
.
% Required packages: amsmath, amssymb, bm, booktabs.
% Status: Experiment 1 is implemented and running; Experiments 2 and 3 are
% prospective protocols. Replace the explicit status statements with measured
% results once each continuation gate has been passed.

\section{Experiments}
\label{sec:experiments}

We evaluate whether sparse predictive structure can identify time-invariant physical
parameters while progressively removing access to the ground-truth state.  We use a
three-experiment ladder.  Experiment~1 gives both branches the true physical state;
Experiment~2 replaces the predictor-side state by a causal visual encoding while retaining a
true-state target; and Experiment~3 jointly learns the context and single-frame target
representations from pixels.  The physical system, parameter distribution, SPARTAN predictor,
and evaluation logic are held fixed as far as the change in observation contract permits.
Each experiment is run only after its predecessor passes a predeclared continuation gate.  This
staging separates failures of causal parameter identification from failures of perception,
tracking, or jointly learned target geometry.

At the time of writing, Experiment~1 is implemented and its attention-logit sweep is running.
The later experiments below specify the intended confirmatory protocols and must not be read as
completed empirical results.

\subsection{Bounce}
\label{sec:experiments:bounce}

\subsubsection{Data}
\label{sec:experiments:bounce:data}

An episode contains $N=5$ elastic balls in the unit square.  We use $X_t$ for a rendered
frame, $Z_t=(z_t^1,\ldots,z_t^N)$ for the true kinematic state, and
$\boldsymbol{m}=(m_1,\ldots,m_N)$ for the time-invariant physical parameters, where
\begin{equation}
    X_t\in[0,1]^{3\times64\times64},
    \qquad
    z_t^i=(p_{x,t}^i,p_{y,t}^i,v_{x,t}^i,v_{y,t}^i)\in\mathbb{R}^{4},
    \qquad
    m_i>0.
    \label{eq:bounce-observations}
\end{equation}
For every episode and object, the implemented Baumgartner-aligned configuration samples
\begin{equation}
    \epsilon_i\sim\mathcal{N}(0,1),
    \qquad
    m_i=\operatorname{clip}_{[0.5,3.0]}(1.5+0.5\epsilon_i),
    \qquad
    r_i=r_0\frac{m_i}{m_{\mathrm{ref}}},
    \label{eq:bounce-mass-radius}
\end{equation}
with $r_0=0.08$ and the episode-independent reference $m_{\mathrm{ref}}=1.5$.
Equation~\eqref{eq:bounce-mass-radius} is a clipped Gaussian, rather than a truncated
Gaussian, and therefore places atoms at the clipping boundaries.  The numerical mass
distribution is our declared choice because its exact parameters are not specified by the
source Bounce experiment.

The initial velocity has fixed norm $0.7$ and uniformly random direction,
\begin{equation}
    \varphi_i\sim\mathcal{U}[0,2\pi),
    \qquad
    v_0^i=0.7(\cos\varphi_i,\sin\varphi_i).
\end{equation}
The current generator sequentially samples an initial centre uniformly from
$[1.05r_i,1-1.05r_i]^2$ and accepts it only if
$\lVert p_0^i-p_0^j\rVert_2>1.1(r_i+r_j)$ for all previously placed objects $j$.
Consequently, the support of the initial positions is itself mass-dependent.  We report this
fact rather than describing the initial state as independent of mass.  To isolate evidence from
the subsequent dynamics, the final study also includes a mass-independent-initialization
control: all wall margins and pair-separation constraints are computed using the common
worst-case radius $r_{\max}=r_0(3.0/1.5)$, irrespective of the sampled masses.  This control is
required before attributing recovery specifically to transition dynamics rather than to the
initial-state support.

The simulator records $T=60$ states at intervals $\Delta t=0.1$, with eight physics
substeps per recorded transition.  Free motion is followed by exact reflection of wall
overshoot about $r_i$ or $1-r_i$.  For a colliding pair, let
\begin{equation}
    n_{ij}=\frac{p_i-p_j}{\lVert p_i-p_j\rVert_2},
    \qquad
    u_{ij}=(v_i-v_j)^\top n_{ij}<0.
\end{equation}
The restitution-one velocity update is
\begin{align}
    v_i' &= v_i-\frac{2m_j}{m_i+m_j}u_{ij}n_{ij},
    &
    v_j' &= v_j+\frac{2m_i}{m_i+m_j}u_{ij}n_{ij}.
    \label{eq:elastic-collision}
\end{align}
Before applying this impulse, any penetration is projected back to the contact surface
$\lVert p_i-p_j\rVert_2=r_i+r_j$ using inverse-mass weights.  The corrected simulator is
versioned so that pre-generated trajectories from older collision rules cannot be mixed with
the current data.

Each episode supplies
\begin{equation}
    (X_{0:T-1},Z_{0:T-1},\boldsymbol{m},C_{0:T-2}),
    \qquad
    C_t\in\{0,1\}^{N\times N},
\end{equation}
where $C_{t,ij}=1$ records a pair contact during transition $t\rightarrow t+1$.
Off-diagonal entries are symmetric.  In the mass-dependent-radius environment,
$C_{t,ii}=1$ records a wall contact because its location depends on $m_i$ through $r_i$.
We use the first $T_{\mathrm{par}}=30$ observations to infer one parameter representation and
evaluate $K=T-T_{\mathrm{par}}=30$ transitions,
\begin{equation}
    \mathcal{I}=\{T_{\mathrm{par}}-1,\ldots,T-2\}=\{29,\ldots,58\}.
    \label{eq:transition-index-set}
\end{equation}
Thus, the parameter encoder sees observations $0,\ldots,29$, while the supervised pairs in
Experiment~1 are $(Z_{29},Z_{30}),\ldots,(Z_{58},Z_{59})$.  These are 30
teacher-forced one-step predictions, not one 30-step open-loop rollout.

\paragraph{Why absolute masses are observable.}
If all physical radii were equal, multiplying every mass by a common positive constant would
leave Equation~\eqref{eq:elastic-collision} unchanged; elastic impulses would then identify
mass ratios but not absolute scale.  The implemented relation
$r_i=r_0m_i/m_{\mathrm{ref}}$ removes that symmetry because a common mass rescaling also
changes pair-contact distances, wall-contact locations, penetration correction, and the support
of initial positions.  The current state experiment sets rendering off and never supplies a
literal radius coordinate, but its observed positions and velocities are nevertheless generated
by this mass-dependent geometry.  We therefore claim identification in the
mass-proportional-physical-radius environment, not in the harder equal-radius environment.

For Experiments~2 and~3, physical and rendered radii must be separated.  The simulator retains
Equation~\eqref{eq:bounce-mass-radius}, while every object is drawn with the same
mass-independent glyph and rendered radius.  The primary visual data must also contain no
dataset-global visual signature tied to simulator row index.  Hence a single frame cannot reveal
mass from the object's glyph, although a sequence can reveal mass through centre motion and
hidden contact geometry.  A visible-radius version and an equal-physical-radius version are
reported only as labelled controls.  The current data implementation uses a single radius switch
for physics and rendering; splitting those two settings and versioning the resulting data are
prerequisites for the visual experiments.

\subsubsection{Ground-truth local graphs}
\label{sec:experiments:bounce:graphs}

The destination row is indexed first in every graph.  From $C_t$, we define the true
state-to-state and mass-to-state local graphs as
\begin{align}
    G^{Z\rightarrow Z}_{t,ij}
        &= \mathbf{1}\{i=j\ \lor\ C_{t,ij}=1\},
        \\
    G^{m\rightarrow Z}_{t,ij}
        &= \mathbf{1}\!\left\{C_{t,ij}=1\ \lor\
             \left(i=j\ \land\ \bigvee_k C_{t,ik}=1\right)\right\}.
    \label{eq:bounce-ground-truth-graphs}
\end{align}
A pair collision therefore gives each affected state an own-mass and partner-mass parent; a
mass-dependent wall collision gives the affected state an own-mass parent.  Self-state edges are
always present because free motion depends on the object's previous state.

\subsubsection{Shared model comparisons and metrics}
\label{sec:experiments:shared-protocol}

Every scientific comparison uses the same three predictor modes under matched initialization
rules, training budgets, and data splits:
\begin{enumerate}
    \item \emph{Dense}: $A^{(\ell)}\equiv\boldsymbol{1}$ in every SPARTAN layer and no
          path penalty.
    \item \emph{Token-local}: $A^{(\ell)}\equiv\boldsymbol{0}$, so only each token's
          residual--MLP path remains.  This reference has no access to other state or parameter
          tokens, but its own state may contain statistical traces of earlier mass-dependent
          geometry; it is therefore more precise than calling the reference fully mass-blind.
    \item \emph{Sparse}: learned hard Bernoulli gates with the SPARTAN path objective and a
          GECO-style constraint controller.
\end{enumerate}

There are $N$ state tokens and $N$ parameter tokens.  For $L_{\mathrm{sp}}=3$ SPARTAN
layers, the multilayer reachability matrix is
\begin{equation}
    \overline A
    =\bigl(A^{(L_{\mathrm{sp}})}+I\bigr)\cdots
     \bigl(A^{(1)}+I\bigr),
    \label{eq:spartan-path-matrix}
\end{equation}
where $\overline A_{ij}>0$ means that at least one permitted computational path connects source
token $j$ to decoded output $i$.  The optimized path count includes only the $N$ decoded state
rows,
\begin{equation}
    \mathcal{L}_{\mathrm{path}}
    =\mathbb{E}\!\left[
       \sum_{i=1}^{N}\sum_{j=1}^{2N}\overline A_{ij}
     \right],
    \qquad
    \rho_{\mathrm{path}}
    =\frac{1}{2N^2}\sum_{i=1}^{N}\sum_{j=1}^{2N}
       \mathbf{1}\{\overline A_{ij}>0\}.
    \label{eq:path-objective-density}
\end{equation}
The second quantity discards path multiplicity and is the primary pruning metric.  With ten
tokens, the token-local reference has $\overline A=I_{10}$ and
$\rho_{\mathrm{path}}=5/50=0.1$, while the dense model has
$\rho_{\mathrm{path}}=1$.  The corresponding unnormalised path objectives are 5 and 6655:
if $J$ is the $10\times10$ all-ones matrix, then
$(J+I)^3=I+133J$, each decoded row sums to 1331, and five decoded rows sum to 6655.
Reachability over all output rows is retained as a diagnostic only, since parameter-token outputs
are not decoded.

For a fixed target space, training minimizes
\begin{equation}
    \mathcal{L}
    =\mathcal{L}_{\mathrm{pred}}
     +\lambda_{\mathrm{logit}}\mathcal{L}_{\mathrm{logit}}
     +\lambda_{\mathrm{reg}}\mathcal{R}
     +\Lambda^{-1}\mathcal{L}_{\mathrm{path}},
    \label{eq:shared-training-objective}
\end{equation}
when sparsity is enabled.  Here $\mathcal{R}$ is VISReg on learned
representations.  Fixed raw-state targets are not regularized: applying VISReg to them produces
no model gradient and only adds a stochastic constant.  The attention-logit penalty averages
$\exp(a)+\exp(-a)$ over layer logits $a$ and has theoretical minimum 2; the implementation uses
a linear continuation beyond $|a|=30$ for numerical stability.

The dual constraint excludes both VISReg and the path penalty,
\begin{equation}
    c=\mathcal{L}_{\mathrm{pred}}
      +\lambda_{\mathrm{logit}}\mathcal{L}_{\mathrm{logit}}
      \leq \tau.
    \label{eq:raw-geco-constraint}
\end{equation}
For every architecture and seed, $\tau$ is recalibrated as the held-out constraint of a
converged dense reference, and the sparse run is launched only if the dense constraint is below
the token-local constraint under a predeclared paired uncertainty test.  The dual variable starts
at $\Lambda=10^6$, making the initial path weight $10^{-6}$, and increases or decreases according
to a moving average of $c-\tau$.  A Boolean indicator that sparsity is enabled is not evidence
that pruning occurred.

The coefficient $\lambda_{\mathrm{logit}}$ is selected without mass labels.  The dense sweep
contains an exact zero control, rejects candidates whose prediction error is more than 5\% worse
than that control, and selects the smallest Pareto coefficient attaining 90\% of the best
admissible reduction in $\mathcal{L}_{\mathrm{logit}}-2$.  We repeat or revalidate this screen
whenever the token geometry, matching rule, or learned target space changes.  Neither masses nor
graph labels select hyperparameters, checkpoints, or seeds.

\paragraph{Parameter recovery.}
In Experiment~1, the learned parameter vector consists of five persistent global coordinates
$\widehat{\boldsymbol\theta}\in\mathbb{R}^{5}$.  On a probe-training fold we fit a scalar
one-hidden-layer MLP for every pair $(\widehat\theta_j,m_i)$.  A disjoint alignment fold yields
a nonlinear $R^2$ matrix and one global Hungarian bijection $\sigma$; the frozen bijection is
scored on a third fold:
\begin{equation}
    \operatorname{MCC}_{1:1}
      =\frac{1}{N}\sum_{i=1}^{N}R^2_{\sigma(i),i}.
    \label{eq:strict-mcc}
\end{equation}
Negative $R^2$ values are clipped to zero.  This is stricter than a mean--maximum score because
one learned coordinate cannot explain several masses and the permutation cannot change between
episodes.  The same frozen $\sigma$ aligns parameter columns before computing graph error.  In
the track-attached visual experiments, the single trajectory assignment instead pairs each
learned track scalar with its physical object; the nonlinear probe is then scored over held-out
episode--track pairs.

We report raw prediction error, the exact constraint in
Equation~\eqref{eq:raw-geco-constraint}, $\rho_{\mathrm{path}}$, state and parameter structural
Hamming distance (SHD), all five assigned nonlinear $R^2$ values, and
$\operatorname{MCC}_{1:1}$.  SHD is the mean unnormalised number of mismatched entries in the
$N\times N$ local graph and therefore lies in $[0,25]$.  Decodability alone does not establish
causal use.  We consequently evaluate three interventions: permuting learned parameter vectors
between episodes, replacing them by their training mean, and disabling retained
parameter-to-state paths.  Effects are reported separately on contact and non-contact
transitions.

Seed 0 is used for development.  Once a protocol is frozen, eight fresh paired seeds (1--8) are
reported, including valid optimization failures.  Dense, token-local, and sparse checkpoints are
evaluated on the same final 5,000-episode split.  Paired differences use a percentile bootstrap
over seed pairs with 10,000 resamples and a fixed analysis seed.  Interrupted jobs, corrupt
artifacts, or provenance mismatches are rerun; reproducible optimization failures remain in the
denominator.

\begin{table}[t]
    \centering
    \small
    \begin{tabular}{p{0.18\linewidth}p{0.22\linewidth}p{0.23\linewidth}p{0.27\linewidth}}
        \toprule
        Experiment & Predictor-side observation & Prediction target & Strongest supported claim \\
        \midrule
        1: true state & $Z_{0:t}$ & true $Z_{t+1}$ & Identification and sparse use conditional on ordered oracle state tracks. \\
        2: visual context & $X_{0:t}$ & true $Z_{t+1}$ & Visual state estimation and mass discovery with state-supervised training. \\
        3: fully visual & $X_{0:t}$ & $E_\phi(X_{t+1})$ & Jointly learned visual mass representation and sparse causal use, up to trajectory permutation. \\
        \bottomrule
    \end{tabular}
    \caption{The three experiments change the observation contract one step at a time.}
    \label{tab:experiment-ladder}
\end{table}

\subsection{Experiment 1: true-state context and true-state target}
\label{sec:experiments:true-state}

\paragraph{Question and source of mass information.}
This experiment asks whether the parameter encoder and SPARTAN can recover and use the masses
when perception, state grounding, and slot correspondence are given.  The model receives only
$[p_x,p_y,v_x,v_y]$ rows, but those trajectories contain mass information through collision
impulses, mass-dependent contact and wall geometry, and---in the current primary generator---the
mass-dependent initial-position support.  The mass-independent-initialization control isolates
the latter contribution.

\paragraph{Model.}
A shared linear map embeds each context state into 32 dimensions,
\begin{equation}
    e_t^i=W_c z_t^i+b_c\in\mathbb{R}^{32}.
\end{equation}
Five learned latent queries attend to all $T_{\mathrm{par}}N=150$ context tokens after learned
temporal and source-track embeddings, and a scalar head produces
\begin{equation}
    \widehat{S}^{\mathrm{ph}}
    \equiv\widehat{\boldsymbol\theta}
    =P_\eta(e_{0:T_{\mathrm{par}}-1})\in\mathbb{R}^{5\times1}.
    \label{eq:experiment-one-parameters}
\end{equation}
These are persistent global coordinates; coordinate $j$ is not declared to be the mass of input
row $j$.  Independent learned addresses distinguish the five state nodes and five parameter
nodes without privileging any of the 25 parameter-to-state relations.  SPARTAN separately
projects four-dimensional state tokens and scalar parameter tokens into a 512-dimensional
workspace, applies three sparse-attention layers with three-linear-layer 512-wide MLPs, and
projects the five decoded state tokens back to $\mathbb{R}^{4}$:
\begin{equation}
    \widehat Z_{t+1}
      =f_\gamma(Z_t,\widehat{\boldsymbol\theta}),
    \qquad t\in\mathcal I.
    \label{eq:experiment-one-predictor}
\end{equation}
The same $\widehat{\boldsymbol\theta}$ is reused for all 30 transitions in an episode.

Training uses aligned raw-state MSE,
\begin{equation}
    \mathcal{L}^{(1)}_{\mathrm{pred}}
    =\frac{1}{|\mathcal I|N\,4}
      \sum_{t\in\mathcal I}\sum_{i=1}^{N}
      \lVert\widehat z_{t+1}^i-z_{t+1}^i\rVert_2^2,
    \label{eq:experiment-one-loss}
\end{equation}
with every prediction anchored at the corresponding true $Z_t$.  The operative run uses 100,000
training episodes, batch size 16 (480 supervised transitions per update), Adam with learning rate
$5\times10^{-5}$, and 300,000 updates.  No target encoder is present.  VISReg is applied only to
learned context-side representations in the confirmatory protocol.

\paragraph{Current status.}
The state model, global parameter coordinates, dense and token-local references, learned hard
gates, raw constraint, one-to-one mass recovery, and graph evaluation are implemented.  Before a
paper claim, all three checkpoints must be evaluated on the same final split; evaluation artifacts
must be split-specific; paired bootstrap reporting and parameter interventions must be added; and
the gradient-free target VISReg constant must be removed from this regime.

\paragraph{Continuation gate.}
We begin implementation of Experiment~2 only after a complete development pilot shows that
(i) the dense constraint is reliably below the token-local constraint; (ii) the sparse constraint
crosses its freshly calibrated $\tau$; (iii) $\Lambda$ leaves its ceiling downwards and held-out
$\rho_{\mathrm{path}}$ falls for at least three consecutive evaluations; (iv) final density is
strictly between the token-local and dense endpoints; (v) five distinct masses are recovered
under one frozen global permutation; (vi) aligned parameter SHD improves; and (vii) at least one
parameter intervention worsens contact-sensitive prediction.  A failed but understood pilot may
motivate a protocol revision; adding pixels cannot repair an unexplained failure of this oracle
state experiment.

The confirmatory Experiment~1 claim additionally uses the following preregistered targets:
\begin{itemize}
    \item at least six of eight sparse seeds show sustained pruning and finish with the median of
          their last three density evaluations in $(0.10,0.90]$;
    \item the ratio of mean sparse to paired-dense test MSE is at most $1.05$, and at least six of
          eight seeds have a last-three-validation median constraint no larger than $1.05\tau$;
    \item median sparse $\operatorname{MCC}_{1:1}$ is at least $0.80$, with at least six of eight
          seeds at or above $0.70$;
    \item the paired 95\% interval for sparse minus token-local MCC lies above zero, and mean
          sparse parameter SHD is below both matched references;
    \item a positive paired sparse-minus-dense MCC interval is the strongest result, but a
          structural result remains supported when that interval overlaps zero if the ratio of
          mean sparse to dense MCC is at least $0.95$ and all prediction, pruning, graph, and
          intervention gates pass.
\end{itemize}
The Experiment~1 claim is conditional on ordered oracle kinematic tracks.  It is not yet a visual
or anonymous-object result, and high MCC without an intervention effect is described as parameter
information rather than demonstrated causal use.

\subsection{Permutation-safe readiness checks}
\label{sec:experiments:readiness}

The fixed global coordinates and node addresses in Experiment~1 are appropriate for ordered
simulator rows but not for anonymous visual slots.  Before introducing pixels, we validate the
visual coordinate contract on consistently permuted true-state trajectories.  The visual models
emit one scalar parameter per inferred track.  The state token and parameter token of track $i$
receive the same ephemeral key $\kappa_i$; $\kappa_i$ carries no state or physical information
and permutes with the complete track.  This binding is architectural, whereas the parameter value
and the SPARTAN paths that use it remain learned.

For a trajectory-level raw-state assignment, define the standardized cost
\begin{equation}
    D_e(i,j)=\frac{1}{|\mathcal I|\,4}
      \sum_{t\in\mathcal I}\sum_{a=1}^{4}
      \left(
        \frac{\widehat z_{e,t+1,i,a}-z_{e,t+1,j,a}}{\sigma_a}
      \right)^2,
    \qquad
    \pi_e^*=\underset{\pi\in\mathfrak S_N}{\operatorname{argmin}}
      \sum_{i=1}^{N}D_e(i,\pi(i)).
    \label{eq:trajectory-hungarian-raw}
\end{equation}
The discrete assignment is detached, but gradients pass through the selected predictions.  One
$\pi_e^*$ is held fixed for the complete trajectory and is reused for state loss, parameter
recovery, and every graph axis; per-frame rematching is forbidden.  The readiness suite requires
invariance to a whole-trajectory permutation, a positive penalty for a mid-trajectory switch,
joint equivariance of state tokens, parameter tokens and keys, and correct graph-axis alignment.
We retain exactly five slots until an explicit null-slot contract is introduced.  These are
implementation-validity checks, not a fourth scientific experiment.

\subsection{Experiment 2: visual context and true-state target}
\label{sec:experiments:visual-state}

\paragraph{Question and source of mass information.}
Experiment~2 asks whether a visual encoder can replace every predictor-side true state while the
prediction target remains the fixed physical state.  Because rendered glyphs have equal size and
carry no persistent simulator-row signature, mass must be inferred from a 30-frame history of
motion and contact geometry rather than from a single object's appearance.  The unchanged raw
target makes failures attributable to perception or parameter inference rather than to a moving
target representation.

\paragraph{Model.}
A causal recurrent encoder $Q_\psi$ maps frames to five 32-dimensional tracks,
\begin{equation}
    \widetilde S_{0:t}=Q_\psi(X_{0:t}),
    \qquad
    \widetilde S_t\in\mathbb{R}^{N\times32}.
\end{equation}
A shared grounded head gives the current physical-state estimate, while track-aware temporal and
cross-object pooling gives one scalar parameter per track:
\begin{align}
    Z_t^{\mathrm{ctx}} &=g_\omega(\widetilde S_t)\in\mathbb{R}^{N\times4},
    \\
    \widehat{S}^{\mathrm{ph}}
       &=P_\eta(\widetilde S_{0:T_{\mathrm{par}}-1})
         \in\mathbb{R}^{N\times1}.
    \label{eq:experiment-two-heads}
\end{align}
Only frames $0,\ldots,29$ may affect $\widehat S^{\mathrm{ph}}$.  Later source frames are
available for teacher-forcing $Z_t^{\mathrm{ctx}}$, but perturbing them must leave the parameter
estimate unchanged.  With the shared ephemeral track keys from
Section~\ref{sec:experiments:readiness}, SPARTAN uses a four-dimensional state interface, a scalar
parameter interface, and a 512-dimensional internal workspace:
\begin{equation}
    \widehat Z_{t+1}
      =f_\gamma(Z_t^{\mathrm{ctx}},\widehat S^{\mathrm{ph}})
      \in\mathbb{R}^{N\times4}.
    \label{eq:experiment-two-prediction}
\end{equation}
There is no target encoder in this experiment.  The target is the unchanged $Z_{t+1}$, and
training remains one-step teacher-forced; longer open-loop rollouts are evaluation diagnostics.

The trajectory assignment in Equation~\eqref{eq:trajectory-hungarian-raw} is selected once per
episode and reused for both current-state grounding and next-state prediction.  With targets
gathered under the detached assignment, the objective is
\begin{equation}
    \mathcal{L}^{(2)}
      =\mathcal{L}^{(2)}_{\mathrm{pred}}
       +\lambda_{\mathrm{ground}}\mathcal{L}_{\mathrm{ground}}
       +\lambda_{\mathrm{logit}}\mathcal{L}_{\mathrm{logit}}
       +\lambda_{\mathrm{reg}}\mathcal{R}(\widetilde S)
       +\Lambda^{-1}\mathcal{L}_{\mathrm{path}}.
    \label{eq:experiment-two-objective}
\end{equation}
The GECO constraint contains only raw next-state MSE and the weighted logit penalty, as in
Equation~\eqref{eq:raw-geco-constraint}; grounding and VISReg remain outside it.  There is no
target-side VISReg because the target is fixed data.

\paragraph{Continuation gate.}
Experiment~2 must satisfy all shared dense, token-local, sparse, pruning, graph, and intervention
checks under a freshly selected $\lambda_{\mathrm{logit}}$ and freshly calibrated $\tau$.  In
addition, we require: near-complete five-object coverage and a low trajectory switch rate,
including at contacts; accurate grounded position and velocity without target-frame leakage;
stronger held-out mass recovery from $\widehat S^{\mathrm{ph}}$ than incremental mass information
in the four-dimensional state bottleneck; no change in $\widehat S^{\mathrm{ph}}$ when frames after
$T_{\mathrm{par}}$ are perturbed; and success with mass-dependent physical radii but
mass-independent uniform rendering.  Tracking thresholds and grounding tolerances are frozen on
the development split before paired confirmatory seeds.

Only after a complete pilot meets this conjunction do we proceed to the fully visual target.
The supported claim is that frame inputs suffice at inference for state estimation and sparse mass
discovery under state-supervised training.  It is not a self-supervised result: true states define
the grounding loss, prediction target, and training assignment.  Removing the current-state
grounding is retained as an ablation, not as a separate main experiment.

\subsection{Optional diagnostic: learned target of the true future state}
\label{sec:experiments:state-target-diagnostic}

If Experiment~3 fails, we insert a small diagnostic between Experiments~2 and~3.  A jointly
learned shared per-object encoder maps the true $Z_{t+1}\in\mathbb{R}^{N\times4}$ to a
32-dimensional target, while the source remains visual.  Dense and token-local runs test whether
a simultaneously moving target can remain non-collapsed and comparable when object rows are
already given.  If this diagnostic succeeds but the image target fails, slot formation or target
tracking is implicated; if both fail, the moving target geometry or calibration is implicated.
Because this diagnostic still trains on true states and does not test single-frame visual slot
formation, it is not a mandatory experiment and does not automatically receive a full sparse
multi-seed study.

\subsection{Experiment 3: jointly learned visual context and target}
\label{sec:experiments:fully-visual}

\paragraph{Question.}
The final experiment asks whether mass recovery, parameter use, and SPARTAN pruning survive when
all trainable components are learned jointly from frames.  During training, true states, masses,
contacts, renderer masks, and simulator identities are unavailable to the model, loss, and
matching procedure.  They are used only for held-out evaluation.  The long-context encoder, the
single-frame observation/target encoder, parameter encoder, temporal lift, and SPARTAN are all
optimized simultaneously; an EMA or frozen target is a labelled control rather than the primary
model.

\paragraph{Independent single-frame observations and generic target tracks.}
The observation encoder is called independently on one image at a time,
\begin{equation}
    \{(o_t^r,M_t^r)\}_{r=1}^{N}=E_\phi(X_t),
    \qquad
    o_t^r\in\mathbb{R}^{32},
    \label{eq:single-frame-encoder}
\end{equation}
where $M_t^r$ is the slot assignment map.  In particular, the target call for transition
$t\rightarrow t+1$ receives only $X_{t+1}$.  The loss may associate outputs from several such
independent calls, but the encoder itself never receives a target sequence or context history.

Let $\mu(M)$ denote the spatial centroid of an assignment map.  Since the currently normalized
map sums to one for every slot, its total mass is not a valid presence score; any future
confidence statistic must instead be derived from pre-normalization assignment mass, logits, or
entropy.  Target tracks are assembled sequentially with a detached Hungarian assignment
\begin{equation}
    \rho_{e,t}^*
      =\underset{\rho\in\mathfrak S_N}{\operatorname{argmin}}
        \sum_{i=1}^{N}\mathcal{C}_{e,t}(i,\rho(i)),
    \qquad
    \overline o_{e,t}^{\,i}=o_{e,t}^{\rho_{e,t}^*(i)}.
    \label{eq:target-track-assignment}
\end{equation}
For the primary Bounce experiment, $\mathcal C$ is geometry-only: it combines the distance
between the candidate centroid and the constant-velocity prediction
$2\mu(M_{t-1}^i)-\mu(M_{t-2}^i)$ with overlap between the candidate support and a correspondingly
translated previous support.  The first two frames use the analogous position-only costs.  A
normalized learned-feature distance is a possible extension for more complex scenes, but it is
not a fixed renderer-identity key and must generalize under held-out appearance changes.  The
primary Bounce renderer uses uniform object glyphs, so no hand-specified appearance category can
solve Equation~\eqref{eq:target-track-assignment}.

\paragraph{A type-closed dynamical state.}
A single frame supplies position-like evidence but cannot determine velocity.  We therefore use a
shared two-frame lift,
\begin{equation}
    S_t^i=L_\xi\!\left([\nu(\overline o_{t-1}^{\,i}),
                         \nu(\overline o_t^{\,i})]\right)
       \in\mathbb{R}^{32},
    \label{eq:temporal-lift}
\end{equation}
where $\nu$ is a predeclared per-slot normalization.  In parallel, a separate causal recurrent
context encoder and track-aware parameter pooler use only the first 30 frames,
\begin{equation}
    \widetilde S_{0:T_{\mathrm{par}}-1}
       =Q_\psi(X_{0:T_{\mathrm{par}}-1}),
    \qquad
    \widehat S^{\mathrm{ph}}
       =P_\eta(\widetilde S_{0:T_{\mathrm{par}}-1})
       \in\mathbb{R}^{N\times1}.
    \label{eq:visual-parameter-context}
\end{equation}
The parameter and dynamical-state tokens use the same permutation-carrying track keys described
in Section~\ref{sec:experiments:readiness}.  SPARTAN works internally at width 512 and predicts
the next 32-dimensional observation slots,
\begin{equation}
    \widehat O_{t+1}=f_\gamma(S_t,\widehat S^{\mathrm{ph}})
       \in\mathbb{R}^{N\times32}.
    \label{eq:fully-visual-prediction}
\end{equation}
An autoregressive prediction is fed back through the same lift,
\begin{equation}
    \widehat S_{t+1}^{\,i}
       =L_\xi\!\left([\nu(\overline o_t^{\,i}),
                       \nu(\widehat o_{t+1}^{\,i})]\right),
    \qquad
    \widehat O_{t+2}=f_\gamma(\widehat S_{t+1},\widehat S^{\mathrm{ph}}).
    \label{eq:type-closed-rollout}
\end{equation}
This avoids feeding a target-space prediction directly into a context-head state whose equal
width does not guarantee equal meaning.

After independently encoded targets have been assembled into tracks, a second detached
trajectory-level Hungarian assignment $\pi_e^*$ matches predicted tracks to target tracks by
their mean normalized latent distance across all prediction times.  This one assignment is held
fixed for the complete loss and graph alignment.  The within-target assignments $\rho_{e,t}^*$
and the predictor-to-target assignment $\pi_e^*$ are distinct; replacing either with independent
per-frame prediction matching could hide track switches.

\paragraph{Objective and common-start comparison.}
Let $\mathcal{L}_{\mathrm{TF}}$ denote trajectory-matched one-step latent MSE and let
$\mathcal{L}^{(2)}_{\mathrm{AR}}$ denote the two-step loss obtained from
Equation~\eqref{eq:type-closed-rollout}.  The predictive objective is
\begin{equation}
    \mathcal{L}^{(3)}_{\mathrm{pred}}
      =\mathcal{L}_{\mathrm{TF}}
       +\beta\mathcal{L}^{(2)}_{\mathrm{AR}},
    \label{eq:fully-visual-predictive-loss}
\end{equation}
where $\beta$ is fixed before the reference runs.  Longer open-loop horizons through multiple
contacts are evaluated but are not initially backpropagated end to end.  VISReg is computed per
aligned track across episodes and time for both learned branches; flattening all slot addresses
together would not exclude a static slot-codebook collapse.

For the learned target, the dual uses a detached within-track variance scale,
\begin{align}
    V_{\mathrm{within}}
      &=\frac{1}{N\,32}\sum_{i=1}^{N}\sum_{a=1}^{32}
        \operatorname{Var}_{e,t}(\overline o_{e,t,i,a}),
        \\
    c_3
      &=\frac{\mathcal{L}^{(3)}_{\mathrm{pred}}}
              {\operatorname{sg}(V_{\mathrm{within}})}
        +\lambda_{\mathrm{logit}}\mathcal{L}_{\mathrm{logit}}
      \leq\tau_3.
    \label{eq:fully-visual-constraint}
\end{align}
This scale normalization does not by itself make two independently learned target geometries
comparable.  We therefore first train all representation modules with a dense SPARTAN from
scratch, then clone one common checkpoint into continued dense, token-local, and sparse legs.
The target and context encoders remain trainable in the primary sparse leg, and $\tau_3$ is
measured in the shared checkpoint geometry.  A frozen-target pruning leg is reported as a
fixed-ruler audit, not as a replacement for the fully joint result.  Throughout training we
measure target drift, within-track variance, effective rank, and fixed held-out probes so that a
smaller loss cannot be misread as equal prediction quality after target collapse or semantic
drift.

\paragraph{Final acceptance gate.}
The fully visual result is accepted only if all of the following hold:
\begin{itemize}
    \item every trainable branch receives gradients, and aligned context and target slots retain
          non-trivial within-track variance, effective rank, and object-position information;
    \item the target encoder is called on one frame at a time, later frames cannot affect
          $\widehat S^{\mathrm{ph}}$, and no true state, mass, contact, mask, or simulator identity
          enters the training assignment or loss;
    \item all five objects are covered and the target and predictor trajectory switch rates are
          negligible outside explicitly flagged observational ambiguities, including at contacts;
    \item position is decodable from one-frame slots, while velocity is substantially more
          decodable from the two-frame state in Equation~\eqref{eq:temporal-lift} than from one
          frame alone;
    \item the dense constraint is reliably below the token-local constraint in the common target
          geometry, and the sparse model reaches its fresh $\tau_3$, prunes materially, and
          remains strictly above the token-local density endpoint;
    \item masses are recoverable from $\widehat S^{\mathrm{ph}}$ but not from a fixed
          single-frame glyph signature, and parameter permutation, mean replacement, and retained
          path ablations worsen contact-sensitive prediction;
    \item open-loop latent error and norm are reported by horizon through multiple contacts, and
          dense, token-local, and sparse target representations remain sufficiently comparable for
          the constrained-loss comparison to be meaningful;
    \item the result reproduces over paired confirmatory seeds and under uniform and held-out
          heterogeneous-rendering controls.
\end{itemize}

For visually indistinguishable objects at a genuinely ambiguous crossing or contact, recovery of
an arbitrary simulator label is not identifiable from the observations.  In such cases we report
set-level mass recovery and ambiguity-aware tracking metrics rather than using a renderer-specific
identity shortcut.  Subject to this qualification, Experiment~3 supports the intended final
claim: jointly learned frame representations contain a track-attached, time-invariant mass
representation that SPARTAN uses through a pruned causal predictor, with true simulator variables
used only to evaluate that claim.
.
% Required packages: amsmath, amssymb, bm, booktabs.
% Status: Experiment 1 is implemented and running; Experiments 2 and 3 are
% prospective protocols. Replace the explicit status statements with measured
% results once each continuation gate has been passed.

\section{Experiments}
\label{sec:experiments}

We evaluate whether sparse predictive structure can identify time-invariant physical
parameters while progressively removing access to the ground-truth state.  We use a
three-experiment ladder.  Experiment~1 gives both branches the true physical state;
Experiment~2 replaces the predictor-side state by a causal visual encoding while retaining a
true-state target; and Experiment~3 jointly learns the context and single-frame target
representations from pixels.  The physical system, parameter distribution, SPARTAN predictor,
and evaluation logic are held fixed as far as the change in observation contract permits.
Each experiment is run only after its predecessor passes a predeclared continuation gate.  This
staging separates failures of causal parameter identification from failures of perception,
tracking, or jointly learned target geometry.

At the time of writing, Experiment~1 is implemented and its attention-logit sweep is running.
The later experiments below specify the intended confirmatory protocols and must not be read as
completed empirical results.

\subsection{Bounce}
\label{sec:experiments:bounce}

\subsubsection{Data}
\label{sec:experiments:bounce:data}

An episode contains $N=5$ elastic balls in the unit square.  We use $X_t$ for a rendered
frame, $Z_t=(z_t^1,\ldots,z_t^N)$ for the true kinematic state, and
$\boldsymbol{m}=(m_1,\ldots,m_N)$ for the time-invariant physical parameters, where
\begin{equation}
    X_t\in[0,1]^{3\times64\times64},
    \qquad
    z_t^i=(p_{x,t}^i,p_{y,t}^i,v_{x,t}^i,v_{y,t}^i)\in\mathbb{R}^{4},
    \qquad
    m_i>0.
    \label{eq:bounce-observations}
\end{equation}
For every episode and object, the implemented Baumgartner-aligned configuration samples
\begin{equation}
    \epsilon_i\sim\mathcal{N}(0,1),
    \qquad
    m_i=\operatorname{clip}_{[0.5,3.0]}(1.5+0.5\epsilon_i),
    \qquad
    r_i=r_0\frac{m_i}{m_{\mathrm{ref}}},
    \label{eq:bounce-mass-radius}
\end{equation}
with $r_0=0.08$ and the episode-independent reference $m_{\mathrm{ref}}=1.5$.
Equation~\eqref{eq:bounce-mass-radius} is a clipped Gaussian, rather than a truncated
Gaussian, and therefore places atoms at the clipping boundaries.  The numerical mass
distribution is our declared choice because its exact parameters are not specified by the
source Bounce experiment.

The initial velocity has fixed norm $0.7$ and uniformly random direction,
\begin{equation}
    \varphi_i\sim\mathcal{U}[0,2\pi),
    \qquad
    v_0^i=0.7(\cos\varphi_i,\sin\varphi_i).
\end{equation}
The current generator sequentially samples an initial centre uniformly from
$[1.05r_i,1-1.05r_i]^2$ and accepts it only if
$\lVert p_0^i-p_0^j\rVert_2>1.1(r_i+r_j)$ for all previously placed objects $j$.
Consequently, the support of the initial positions is itself mass-dependent.  We report this
fact rather than describing the initial state as independent of mass.  To isolate evidence from
the subsequent dynamics, the final study also includes a mass-independent-initialization
control: all wall margins and pair-separation constraints are computed using the common
worst-case radius $r_{\max}=r_0(3.0/1.5)$, irrespective of the sampled masses.  This control is
required before attributing recovery specifically to transition dynamics rather than to the
initial-state support.

The simulator records $T=60$ states at intervals $\Delta t=0.1$, with eight physics
substeps per recorded transition.  Free motion is followed by exact reflection of wall
overshoot about $r_i$ or $1-r_i$.  For a colliding pair, let
\begin{equation}
    n_{ij}=\frac{p_i-p_j}{\lVert p_i-p_j\rVert_2},
    \qquad
    u_{ij}=(v_i-v_j)^\top n_{ij}<0.
\end{equation}
The restitution-one velocity update is
\begin{align}
    v_i' &= v_i-\frac{2m_j}{m_i+m_j}u_{ij}n_{ij},
    &
    v_j' &= v_j+\frac{2m_i}{m_i+m_j}u_{ij}n_{ij}.
    \label{eq:elastic-collision}
\end{align}
Before applying this impulse, any penetration is projected back to the contact surface
$\lVert p_i-p_j\rVert_2=r_i+r_j$ using inverse-mass weights.  The corrected simulator is
versioned so that pre-generated trajectories from older collision rules cannot be mixed with
the current data.

Each episode supplies
\begin{equation}
    (X_{0:T-1},Z_{0:T-1},\boldsymbol{m},C_{0:T-2}),
    \qquad
    C_t\in\{0,1\}^{N\times N},
\end{equation}
where $C_{t,ij}=1$ records a pair contact during transition $t\rightarrow t+1$.
Off-diagonal entries are symmetric.  In the mass-dependent-radius environment,
$C_{t,ii}=1$ records a wall contact because its location depends on $m_i$ through $r_i$.
We use the first $T_{\mathrm{par}}=30$ observations to infer one parameter representation and
evaluate $K=T-T_{\mathrm{par}}=30$ transitions,
\begin{equation}
    \mathcal{I}=\{T_{\mathrm{par}}-1,\ldots,T-2\}=\{29,\ldots,58\}.
    \label{eq:transition-index-set}
\end{equation}
Thus, the parameter encoder sees observations $0,\ldots,29$, while the supervised pairs in
Experiment~1 are $(Z_{29},Z_{30}),\ldots,(Z_{58},Z_{59})$.  These are 30
teacher-forced one-step predictions, not one 30-step open-loop rollout.

\paragraph{Why absolute masses are observable.}
If all physical radii were equal, multiplying every mass by a common positive constant would
leave Equation~\eqref{eq:elastic-collision} unchanged; elastic impulses would then identify
mass ratios but not absolute scale.  The implemented relation
$r_i=r_0m_i/m_{\mathrm{ref}}$ removes that symmetry because a common mass rescaling also
changes pair-contact distances, wall-contact locations, penetration correction, and the support
of initial positions.  The current state experiment sets rendering off and never supplies a
literal radius coordinate, but its observed positions and velocities are nevertheless generated
by this mass-dependent geometry.  We therefore claim identification in the
mass-proportional-physical-radius environment, not in the harder equal-radius environment.

For Experiments~2 and~3, physical and rendered radii must be separated.  The simulator retains
Equation~\eqref{eq:bounce-mass-radius}, while every object is drawn with the same
mass-independent glyph and rendered radius.  The primary visual data must also contain no
dataset-global visual signature tied to simulator row index.  Hence a single frame cannot reveal
mass from the object's glyph, although a sequence can reveal mass through centre motion and
hidden contact geometry.  A visible-radius version and an equal-physical-radius version are
reported only as labelled controls.  The current data implementation uses a single radius switch
for physics and rendering; splitting those two settings and versioning the resulting data are
prerequisites for the visual experiments.

\subsubsection{Ground-truth local graphs}
\label{sec:experiments:bounce:graphs}

The destination row is indexed first in every graph.  From $C_t$, we define the true
state-to-state and mass-to-state local graphs as
\begin{align}
    G^{Z\rightarrow Z}_{t,ij}
        &= \mathbf{1}\{i=j\ \lor\ C_{t,ij}=1\},
        \\
    G^{m\rightarrow Z}_{t,ij}
        &= \mathbf{1}\!\left\{C_{t,ij}=1\ \lor\
             \left(i=j\ \land\ \bigvee_k C_{t,ik}=1\right)\right\}.
    \label{eq:bounce-ground-truth-graphs}
\end{align}
A pair collision therefore gives each affected state an own-mass and partner-mass parent; a
mass-dependent wall collision gives the affected state an own-mass parent.  Self-state edges are
always present because free motion depends on the object's previous state.

\subsubsection{Shared model comparisons and metrics}
\label{sec:experiments:shared-protocol}

Every scientific comparison uses the same three predictor modes under matched initialization
rules, training budgets, and data splits:
\begin{enumerate}
    \item \emph{Dense}: $A^{(\ell)}\equiv\boldsymbol{1}$ in every SPARTAN layer and no
          path penalty.
    \item \emph{Token-local}: $A^{(\ell)}\equiv\boldsymbol{0}$, so only each token's
          residual--MLP path remains.  This reference has no access to other state or parameter
          tokens, but its own state may contain statistical traces of earlier mass-dependent
          geometry; it is therefore more precise than calling the reference fully mass-blind.
    \item \emph{Sparse}: learned hard Bernoulli gates with the SPARTAN path objective and a
          GECO-style constraint controller.
\end{enumerate}

There are $N$ state tokens and $N$ parameter tokens.  For $L_{\mathrm{sp}}=3$ SPARTAN
layers, the multilayer reachability matrix is
\begin{equation}
    \overline A
    =\bigl(A^{(L_{\mathrm{sp}})}+I\bigr)\cdots
     \bigl(A^{(1)}+I\bigr),
    \label{eq:spartan-path-matrix}
\end{equation}
where $\overline A_{ij}>0$ means that at least one permitted computational path connects source
token $j$ to decoded output $i$.  The optimized path count includes only the $N$ decoded state
rows,
\begin{equation}
    \mathcal{L}_{\mathrm{path}}
    =\mathbb{E}\!\left[
       \sum_{i=1}^{N}\sum_{j=1}^{2N}\overline A_{ij}
     \right],
    \qquad
    \rho_{\mathrm{path}}
    =\frac{1}{2N^2}\sum_{i=1}^{N}\sum_{j=1}^{2N}
       \mathbf{1}\{\overline A_{ij}>0\}.
    \label{eq:path-objective-density}
\end{equation}
The second quantity discards path multiplicity and is the primary pruning metric.  With ten
tokens, the token-local reference has $\overline A=I_{10}$ and
$\rho_{\mathrm{path}}=5/50=0.1$, while the dense model has
$\rho_{\mathrm{path}}=1$.  The corresponding unnormalised path objectives are 5 and 6655:
if $J$ is the $10\times10$ all-ones matrix, then
$(J+I)^3=I+133J$, each decoded row sums to 1331, and five decoded rows sum to 6655.
Reachability over all output rows is retained as a diagnostic only, since parameter-token outputs
are not decoded.

For a fixed target space, training minimizes
\begin{equation}
    \mathcal{L}
    =\mathcal{L}_{\mathrm{pred}}
     +\lambda_{\mathrm{logit}}\mathcal{L}_{\mathrm{logit}}
     +\lambda_{\mathrm{reg}}\mathcal{R}
     +\Lambda^{-1}\mathcal{L}_{\mathrm{path}},
    \label{eq:shared-training-objective}
\end{equation}
when sparsity is enabled.  Here $\mathcal{R}$ is VISReg on learned
representations.  Fixed raw-state targets are not regularized: applying VISReg to them produces
no model gradient and only adds a stochastic constant.  The attention-logit penalty averages
$\exp(a)+\exp(-a)$ over layer logits $a$ and has theoretical minimum 2; the implementation uses
a linear continuation beyond $|a|=30$ for numerical stability.

The dual constraint excludes both VISReg and the path penalty,
\begin{equation}
    c=\mathcal{L}_{\mathrm{pred}}
      +\lambda_{\mathrm{logit}}\mathcal{L}_{\mathrm{logit}}
      \leq \tau.
    \label{eq:raw-geco-constraint}
\end{equation}
For every architecture and seed, $\tau$ is recalibrated as the held-out constraint of a
converged dense reference, and the sparse run is launched only if the dense constraint is below
the token-local constraint under a predeclared paired uncertainty test.  The dual variable starts
at $\Lambda=10^6$, making the initial path weight $10^{-6}$, and increases or decreases according
to a moving average of $c-\tau$.  A Boolean indicator that sparsity is enabled is not evidence
that pruning occurred.

The coefficient $\lambda_{\mathrm{logit}}$ is selected without mass labels.  The dense sweep
contains an exact zero control, rejects candidates whose prediction error is more than 5\% worse
than that control, and selects the smallest Pareto coefficient attaining 90\% of the best
admissible reduction in $\mathcal{L}_{\mathrm{logit}}-2$.  We repeat or revalidate this screen
whenever the token geometry, matching rule, or learned target space changes.  Neither masses nor
graph labels select hyperparameters, checkpoints, or seeds.

\paragraph{Parameter recovery.}
In Experiment~1, the learned parameter vector consists of five persistent global coordinates
$\widehat{\boldsymbol\theta}\in\mathbb{R}^{5}$.  On a probe-training fold we fit a scalar
one-hidden-layer MLP for every pair $(\widehat\theta_j,m_i)$.  A disjoint alignment fold yields
a nonlinear $R^2$ matrix and one global Hungarian bijection $\sigma$; the frozen bijection is
scored on a third fold:
\begin{equation}
    \operatorname{MCC}_{1:1}
      =\frac{1}{N}\sum_{i=1}^{N}R^2_{\sigma(i),i}.
    \label{eq:strict-mcc}
\end{equation}
Negative $R^2$ values are clipped to zero.  This is stricter than a mean--maximum score because
one learned coordinate cannot explain several masses and the permutation cannot change between
episodes.  The same frozen $\sigma$ aligns parameter columns before computing graph error.  In
the track-attached visual experiments, the single trajectory assignment instead pairs each
learned track scalar with its physical object; the nonlinear probe is then scored over held-out
episode--track pairs.

We report raw prediction error, the exact constraint in
Equation~\eqref{eq:raw-geco-constraint}, $\rho_{\mathrm{path}}$, state and parameter structural
Hamming distance (SHD), all five assigned nonlinear $R^2$ values, and
$\operatorname{MCC}_{1:1}$.  SHD is the mean unnormalised number of mismatched entries in the
$N\times N$ local graph and therefore lies in $[0,25]$.  Decodability alone does not establish
causal use.  We consequently evaluate three interventions: permuting learned parameter vectors
between episodes, replacing them by their training mean, and disabling retained
parameter-to-state paths.  Effects are reported separately on contact and non-contact
transitions.

Seed 0 is used for development.  Once a protocol is frozen, eight fresh paired seeds (1--8) are
reported, including valid optimization failures.  Dense, token-local, and sparse checkpoints are
evaluated on the same final 5,000-episode split.  Paired differences use a percentile bootstrap
over seed pairs with 10,000 resamples and a fixed analysis seed.  Interrupted jobs, corrupt
artifacts, or provenance mismatches are rerun; reproducible optimization failures remain in the
denominator.

\begin{table}[t]
    \centering
    \small
    \begin{tabular}{p{0.18\linewidth}p{0.22\linewidth}p{0.23\linewidth}p{0.27\linewidth}}
        \toprule
        Experiment & Predictor-side observation & Prediction target & Strongest supported claim \\
        \midrule
        1: true state & $Z_{0:t}$ & true $Z_{t+1}$ & Identification and sparse use conditional on ordered oracle state tracks. \\
        2: visual context & $X_{0:t}$ & true $Z_{t+1}$ & Visual state estimation and mass discovery with state-supervised training. \\
        3: fully visual & $X_{0:t}$ & $E_\phi(X_{t+1})$ & Jointly learned visual mass representation and sparse causal use, up to trajectory permutation. \\
        \bottomrule
    \end{tabular}
    \caption{The three experiments change the observation contract one step at a time.}
    \label{tab:experiment-ladder}
\end{table}

\subsection{Experiment 1: true-state context and true-state target}
\label{sec:experiments:true-state}

\paragraph{Question and source of mass information.}
This experiment asks whether the parameter encoder and SPARTAN can recover and use the masses
when perception, state grounding, and slot correspondence are given.  The model receives only
$[p_x,p_y,v_x,v_y]$ rows, but those trajectories contain mass information through collision
impulses, mass-dependent contact and wall geometry, and---in the current primary generator---the
mass-dependent initial-position support.  The mass-independent-initialization control isolates
the latter contribution.

\paragraph{Model.}
A shared linear map embeds each context state into 32 dimensions,
\begin{equation}
    e_t^i=W_c z_t^i+b_c\in\mathbb{R}^{32}.
\end{equation}
Five learned latent queries attend to all $T_{\mathrm{par}}N=150$ context tokens after learned
temporal and source-track embeddings, and a scalar head produces
\begin{equation}
    \widehat{S}^{\mathrm{ph}}
    \equiv\widehat{\boldsymbol\theta}
    =P_\eta(e_{0:T_{\mathrm{par}}-1})\in\mathbb{R}^{5\times1}.
    \label{eq:experiment-one-parameters}
\end{equation}
These are persistent global coordinates; coordinate $j$ is not declared to be the mass of input
row $j$.  Independent learned addresses distinguish the five state nodes and five parameter
nodes without privileging any of the 25 parameter-to-state relations.  SPARTAN separately
projects four-dimensional state tokens and scalar parameter tokens into a 512-dimensional
workspace, applies three sparse-attention layers with three-linear-layer 512-wide MLPs, and
projects the five decoded state tokens back to $\mathbb{R}^{4}$:
\begin{equation}
    \widehat Z_{t+1}
      =f_\gamma(Z_t,\widehat{\boldsymbol\theta}),
    \qquad t\in\mathcal I.
    \label{eq:experiment-one-predictor}
\end{equation}
The same $\widehat{\boldsymbol\theta}$ is reused for all 30 transitions in an episode.

Training uses aligned raw-state MSE,
\begin{equation}
    \mathcal{L}^{(1)}_{\mathrm{pred}}
    =\frac{1}{|\mathcal I|N\,4}
      \sum_{t\in\mathcal I}\sum_{i=1}^{N}
      \lVert\widehat z_{t+1}^i-z_{t+1}^i\rVert_2^2,
    \label{eq:experiment-one-loss}
\end{equation}
with every prediction anchored at the corresponding true $Z_t$.  The operative run uses 100,000
training episodes, batch size 16 (480 supervised transitions per update), Adam with learning rate
$5\times10^{-5}$, and 300,000 updates.  No target encoder is present.  VISReg is applied only to
learned context-side representations in the confirmatory protocol.

\paragraph{Current status.}
The state model, global parameter coordinates, dense and token-local references, learned hard
gates, raw constraint, one-to-one mass recovery, and graph evaluation are implemented.  Before a
paper claim, all three checkpoints must be evaluated on the same final split; evaluation artifacts
must be split-specific; paired bootstrap reporting and parameter interventions must be added; and
the gradient-free target VISReg constant must be removed from this regime.

\paragraph{Continuation gate.}
We begin implementation of Experiment~2 only after a complete development pilot shows that
(i) the dense constraint is reliably below the token-local constraint; (ii) the sparse constraint
crosses its freshly calibrated $\tau$; (iii) $\Lambda$ leaves its ceiling downwards and held-out
$\rho_{\mathrm{path}}$ falls for at least three consecutive evaluations; (iv) final density is
strictly between the token-local and dense endpoints; (v) five distinct masses are recovered
under one frozen global permutation; (vi) aligned parameter SHD improves; and (vii) at least one
parameter intervention worsens contact-sensitive prediction.  A failed but understood pilot may
motivate a protocol revision; adding pixels cannot repair an unexplained failure of this oracle
state experiment.

The confirmatory Experiment~1 claim additionally uses the following preregistered targets:
\begin{itemize}
    \item at least six of eight sparse seeds show sustained pruning and finish with the median of
          their last three density evaluations in $(0.10,0.90]$;
    \item the ratio of mean sparse to paired-dense test MSE is at most $1.05$, and at least six of
          eight seeds have a last-three-validation median constraint no larger than $1.05\tau$;
    \item median sparse $\operatorname{MCC}_{1:1}$ is at least $0.80$, with at least six of eight
          seeds at or above $0.70$;
    \item the paired 95\% interval for sparse minus token-local MCC lies above zero, and mean
          sparse parameter SHD is below both matched references;
    \item a positive paired sparse-minus-dense MCC interval is the strongest result, but a
          structural result remains supported when that interval overlaps zero if the ratio of
          mean sparse to dense MCC is at least $0.95$ and all prediction, pruning, graph, and
          intervention gates pass.
\end{itemize}
The Experiment~1 claim is conditional on ordered oracle kinematic tracks.  It is not yet a visual
or anonymous-object result, and high MCC without an intervention effect is described as parameter
information rather than demonstrated causal use.

\subsection{Permutation-safe readiness checks}
\label{sec:experiments:readiness}

The fixed global coordinates and node addresses in Experiment~1 are appropriate for ordered
simulator rows but not for anonymous visual slots.  Before introducing pixels, we validate the
visual coordinate contract on consistently permuted true-state trajectories.  The visual models
emit one scalar parameter per inferred track.  The state token and parameter token of track $i$
receive the same ephemeral key $\kappa_i$; $\kappa_i$ carries no state or physical information
and permutes with the complete track.  This binding is architectural, whereas the parameter value
and the SPARTAN paths that use it remain learned.

For a trajectory-level raw-state assignment, define the standardized cost
\begin{equation}
    D_e(i,j)=\frac{1}{|\mathcal I|\,4}
      \sum_{t\in\mathcal I}\sum_{a=1}^{4}
      \left(
        \frac{\widehat z_{e,t+1,i,a}-z_{e,t+1,j,a}}{\sigma_a}
      \right)^2,
    \qquad
    \pi_e^*=\underset{\pi\in\mathfrak S_N}{\operatorname{argmin}}
      \sum_{i=1}^{N}D_e(i,\pi(i)).
    \label{eq:trajectory-hungarian-raw}
\end{equation}
The discrete assignment is detached, but gradients pass through the selected predictions.  One
$\pi_e^*$ is held fixed for the complete trajectory and is reused for state loss, parameter
recovery, and every graph axis; per-frame rematching is forbidden.  The readiness suite requires
invariance to a whole-trajectory permutation, a positive penalty for a mid-trajectory switch,
joint equivariance of state tokens, parameter tokens and keys, and correct graph-axis alignment.
We retain exactly five slots until an explicit null-slot contract is introduced.  These are
implementation-validity checks, not a fourth scientific experiment.

\subsection{Experiment 2: visual context and true-state target}
\label{sec:experiments:visual-state}

\paragraph{Question and source of mass information.}
Experiment~2 asks whether a visual encoder can replace every predictor-side true state while the
prediction target remains the fixed physical state.  Because rendered glyphs have equal size and
carry no persistent simulator-row signature, mass must be inferred from a 30-frame history of
motion and contact geometry rather than from a single object's appearance.  The unchanged raw
target makes failures attributable to perception or parameter inference rather than to a moving
target representation.

\paragraph{Model.}
A causal recurrent encoder $Q_\psi$ maps frames to five 32-dimensional tracks,
\begin{equation}
    \widetilde S_{0:t}=Q_\psi(X_{0:t}),
    \qquad
    \widetilde S_t\in\mathbb{R}^{N\times32}.
\end{equation}
A shared grounded head gives the current physical-state estimate, while track-aware temporal and
cross-object pooling gives one scalar parameter per track:
\begin{align}
    Z_t^{\mathrm{ctx}} &=g_\omega(\widetilde S_t)\in\mathbb{R}^{N\times4},
    \\
    \widehat{S}^{\mathrm{ph}}
       &=P_\eta(\widetilde S_{0:T_{\mathrm{par}}-1})
         \in\mathbb{R}^{N\times1}.
    \label{eq:experiment-two-heads}
\end{align}
Only frames $0,\ldots,29$ may affect $\widehat S^{\mathrm{ph}}$.  Later source frames are
available for teacher-forcing $Z_t^{\mathrm{ctx}}$, but perturbing them must leave the parameter
estimate unchanged.  With the shared ephemeral track keys from
Section~\ref{sec:experiments:readiness}, SPARTAN uses a four-dimensional state interface, a scalar
parameter interface, and a 512-dimensional internal workspace:
\begin{equation}
    \widehat Z_{t+1}
      =f_\gamma(Z_t^{\mathrm{ctx}},\widehat S^{\mathrm{ph}})
      \in\mathbb{R}^{N\times4}.
    \label{eq:experiment-two-prediction}
\end{equation}
There is no target encoder in this experiment.  The target is the unchanged $Z_{t+1}$, and
training remains one-step teacher-forced; longer open-loop rollouts are evaluation diagnostics.

The trajectory assignment in Equation~\eqref{eq:trajectory-hungarian-raw} is selected once per
episode and reused for both current-state grounding and next-state prediction.  With targets
gathered under the detached assignment, the objective is
\begin{equation}
    \mathcal{L}^{(2)}
      =\mathcal{L}^{(2)}_{\mathrm{pred}}
       +\lambda_{\mathrm{ground}}\mathcal{L}_{\mathrm{ground}}
       +\lambda_{\mathrm{logit}}\mathcal{L}_{\mathrm{logit}}
       +\lambda_{\mathrm{reg}}\mathcal{R}(\widetilde S)
       +\Lambda^{-1}\mathcal{L}_{\mathrm{path}}.
    \label{eq:experiment-two-objective}
\end{equation}
The GECO constraint contains only raw next-state MSE and the weighted logit penalty, as in
Equation~\eqref{eq:raw-geco-constraint}; grounding and VISReg remain outside it.  There is no
target-side VISReg because the target is fixed data.

\paragraph{Continuation gate.}
Experiment~2 must satisfy all shared dense, token-local, sparse, pruning, graph, and intervention
checks under a freshly selected $\lambda_{\mathrm{logit}}$ and freshly calibrated $\tau$.  In
addition, we require: near-complete five-object coverage and a low trajectory switch rate,
including at contacts; accurate grounded position and velocity without target-frame leakage;
stronger held-out mass recovery from $\widehat S^{\mathrm{ph}}$ than incremental mass information
in the four-dimensional state bottleneck; no change in $\widehat S^{\mathrm{ph}}$ when frames after
$T_{\mathrm{par}}$ are perturbed; and success with mass-dependent physical radii but
mass-independent uniform rendering.  Tracking thresholds and grounding tolerances are frozen on
the development split before paired confirmatory seeds.

Only after a complete pilot meets this conjunction do we proceed to the fully visual target.
The supported claim is that frame inputs suffice at inference for state estimation and sparse mass
discovery under state-supervised training.  It is not a self-supervised result: true states define
the grounding loss, prediction target, and training assignment.  Removing the current-state
grounding is retained as an ablation, not as a separate main experiment.

\subsection{Optional diagnostic: learned target of the true future state}
\label{sec:experiments:state-target-diagnostic}

If Experiment~3 fails, we insert a small diagnostic between Experiments~2 and~3.  A jointly
learned shared per-object encoder maps the true $Z_{t+1}\in\mathbb{R}^{N\times4}$ to a
32-dimensional target, while the source remains visual.  Dense and token-local runs test whether
a simultaneously moving target can remain non-collapsed and comparable when object rows are
already given.  If this diagnostic succeeds but the image target fails, slot formation or target
tracking is implicated; if both fail, the moving target geometry or calibration is implicated.
Because this diagnostic still trains on true states and does not test single-frame visual slot
formation, it is not a mandatory experiment and does not automatically receive a full sparse
multi-seed study.

\subsection{Experiment 3: jointly learned visual context and target}
\label{sec:experiments:fully-visual}

\paragraph{Question.}
The final experiment asks whether mass recovery, parameter use, and SPARTAN pruning survive when
all trainable components are learned jointly from frames.  During training, true states, masses,
contacts, renderer masks, and simulator identities are unavailable to the model, loss, and
matching procedure.  They are used only for held-out evaluation.  The long-context encoder, the
single-frame observation/target encoder, parameter encoder, temporal lift, and SPARTAN are all
optimized simultaneously; an EMA or frozen target is a labelled control rather than the primary
model.

\paragraph{Independent single-frame observations and generic target tracks.}
The observation encoder is called independently on one image at a time,
\begin{equation}
    \{(o_t^r,M_t^r)\}_{r=1}^{N}=E_\phi(X_t),
    \qquad
    o_t^r\in\mathbb{R}^{32},
    \label{eq:single-frame-encoder}
\end{equation}
where $M_t^r$ is the slot assignment map.  In particular, the target call for transition
$t\rightarrow t+1$ receives only $X_{t+1}$.  The loss may associate outputs from several such
independent calls, but the encoder itself never receives a target sequence or context history.

Let $\mu(M)$ denote the spatial centroid of an assignment map.  Since the currently normalized
map sums to one for every slot, its total mass is not a valid presence score; any future
confidence statistic must instead be derived from pre-normalization assignment mass, logits, or
entropy.  Target tracks are assembled sequentially with a detached Hungarian assignment
\begin{equation}
    \rho_{e,t}^*
      =\underset{\rho\in\mathfrak S_N}{\operatorname{argmin}}
        \sum_{i=1}^{N}\mathcal{C}_{e,t}(i,\rho(i)),
    \qquad
    \overline o_{e,t}^{\,i}=o_{e,t}^{\rho_{e,t}^*(i)}.
    \label{eq:target-track-assignment}
\end{equation}
For the primary Bounce experiment, $\mathcal C$ is geometry-only: it combines the distance
between the candidate centroid and the constant-velocity prediction
$2\mu(M_{t-1}^i)-\mu(M_{t-2}^i)$ with overlap between the candidate support and a correspondingly
translated previous support.  The first two frames use the analogous position-only costs.  A
normalized learned-feature distance is a possible extension for more complex scenes, but it is
not a fixed renderer-identity key and must generalize under held-out appearance changes.  The
primary Bounce renderer uses uniform object glyphs, so no hand-specified appearance category can
solve Equation~\eqref{eq:target-track-assignment}.

\paragraph{A type-closed dynamical state.}
A single frame supplies position-like evidence but cannot determine velocity.  We therefore use a
shared two-frame lift,
\begin{equation}
    S_t^i=L_\xi\!\left([\nu(\overline o_{t-1}^{\,i}),
                         \nu(\overline o_t^{\,i})]\right)
       \in\mathbb{R}^{32},
    \label{eq:temporal-lift}
\end{equation}
where $\nu$ is a predeclared per-slot normalization.  In parallel, a separate causal recurrent
context encoder and track-aware parameter pooler use only the first 30 frames,
\begin{equation}
    \widetilde S_{0:T_{\mathrm{par}}-1}
       =Q_\psi(X_{0:T_{\mathrm{par}}-1}),
    \qquad
    \widehat S^{\mathrm{ph}}
       =P_\eta(\widetilde S_{0:T_{\mathrm{par}}-1})
       \in\mathbb{R}^{N\times1}.
    \label{eq:visual-parameter-context}
\end{equation}
The parameter and dynamical-state tokens use the same permutation-carrying track keys described
in Section~\ref{sec:experiments:readiness}.  SPARTAN works internally at width 512 and predicts
the next 32-dimensional observation slots,
\begin{equation}
    \widehat O_{t+1}=f_\gamma(S_t,\widehat S^{\mathrm{ph}})
       \in\mathbb{R}^{N\times32}.
    \label{eq:fully-visual-prediction}
\end{equation}
An autoregressive prediction is fed back through the same lift,
\begin{equation}
    \widehat S_{t+1}^{\,i}
       =L_\xi\!\left([\nu(\overline o_t^{\,i}),
                       \nu(\widehat o_{t+1}^{\,i})]\right),
    \qquad
    \widehat O_{t+2}=f_\gamma(\widehat S_{t+1},\widehat S^{\mathrm{ph}}).
    \label{eq:type-closed-rollout}
\end{equation}
This avoids feeding a target-space prediction directly into a context-head state whose equal
width does not guarantee equal meaning.

After independently encoded targets have been assembled into tracks, a second detached
trajectory-level Hungarian assignment $\pi_e^*$ matches predicted tracks to target tracks by
their mean normalized latent distance across all prediction times.  This one assignment is held
fixed for the complete loss and graph alignment.  The within-target assignments $\rho_{e,t}^*$
and the predictor-to-target assignment $\pi_e^*$ are distinct; replacing either with independent
per-frame prediction matching could hide track switches.

\paragraph{Objective and common-start comparison.}
Let $\mathcal{L}_{\mathrm{TF}}$ denote trajectory-matched one-step latent MSE and let
$\mathcal{L}^{(2)}_{\mathrm{AR}}$ denote the two-step loss obtained from
Equation~\eqref{eq:type-closed-rollout}.  The predictive objective is
\begin{equation}
    \mathcal{L}^{(3)}_{\mathrm{pred}}
      =\mathcal{L}_{\mathrm{TF}}
       +\beta\mathcal{L}^{(2)}_{\mathrm{AR}},
    \label{eq:fully-visual-predictive-loss}
\end{equation}
where $\beta$ is fixed before the reference runs.  Longer open-loop horizons through multiple
contacts are evaluated but are not initially backpropagated end to end.  VISReg is computed per
aligned track across episodes and time for both learned branches; flattening all slot addresses
together would not exclude a static slot-codebook collapse.

For the learned target, the dual uses a detached within-track variance scale,
\begin{align}
    V_{\mathrm{within}}
      &=\frac{1}{N\,32}\sum_{i=1}^{N}\sum_{a=1}^{32}
        \operatorname{Var}_{e,t}(\overline o_{e,t,i,a}),
        \\
    c_3
      &=\frac{\mathcal{L}^{(3)}_{\mathrm{pred}}}
              {\operatorname{sg}(V_{\mathrm{within}})}
        +\lambda_{\mathrm{logit}}\mathcal{L}_{\mathrm{logit}}
      \leq\tau_3.
    \label{eq:fully-visual-constraint}
\end{align}
This scale normalization does not by itself make two independently learned target geometries
comparable.  We therefore first train all representation modules with a dense SPARTAN from
scratch, then clone one common checkpoint into continued dense, token-local, and sparse legs.
The target and context encoders remain trainable in the primary sparse leg, and $\tau_3$ is
measured in the shared checkpoint geometry.  A frozen-target pruning leg is reported as a
fixed-ruler audit, not as a replacement for the fully joint result.  Throughout training we
measure target drift, within-track variance, effective rank, and fixed held-out probes so that a
smaller loss cannot be misread as equal prediction quality after target collapse or semantic
drift.

\paragraph{Final acceptance gate.}
The fully visual result is accepted only if all of the following hold:
\begin{itemize}
    \item every trainable branch receives gradients, and aligned context and target slots retain
          non-trivial within-track variance, effective rank, and object-position information;
    \item the target encoder is called on one frame at a time, later frames cannot affect
          $\widehat S^{\mathrm{ph}}$, and no true state, mass, contact, mask, or simulator identity
          enters the training assignment or loss;
    \item all five objects are covered and the target and predictor trajectory switch rates are
          negligible outside explicitly flagged observational ambiguities, including at contacts;
    \item position is decodable from one-frame slots, while velocity is substantially more
          decodable from the two-frame state in Equation~\eqref{eq:temporal-lift} than from one
          frame alone;
    \item the dense constraint is reliably below the token-local constraint in the common target
          geometry, and the sparse model reaches its fresh $\tau_3$, prunes materially, and
          remains strictly above the token-local density endpoint;
    \item masses are recoverable from $\widehat S^{\mathrm{ph}}$ but not from a fixed
          single-frame glyph signature, and parameter permutation, mean replacement, and retained
          path ablations worsen contact-sensitive prediction;
    \item open-loop latent error and norm are reported by horizon through multiple contacts, and
          dense, token-local, and sparse target representations remain sufficiently comparable for
          the constrained-loss comparison to be meaningful;
    \item the result reproduces over paired confirmatory seeds and under uniform and held-out
          heterogeneous-rendering controls.
\end{itemize}

For visually indistinguishable objects at a genuinely ambiguous crossing or contact, recovery of
an arbitrary simulator label is not identifiable from the observations.  In such cases we report
set-level mass recovery and ambiguity-aware tracking metrics rather than using a renderer-specific
identity shortcut.  Subject to this qualification, Experiment~3 supports the intended final
claim: jointly learned frame representations contain a track-attached, time-invariant mass
representation that SPARTAN uses through a pruned causal predictor, with true simulator variables
used only to evaluate that claim.
